\documentclass[10pt,letterpaper]{article}
\usepackage[margin=.75in]{geometry}
\usepackage[T1]{fontenc}
\usepackage{lmodern}
\usepackage{xcolor}
\usepackage[hyphens]{url}
\usepackage[colorlinks=true,urlcolor=blue,linkcolor=blue]{hyperref}
\urlstyle{same}
\Urlmuskip=0mu plus 1mu
\setlength{\parindent}{0pt}
\setlength{\parskip}{.45em}
\emergencystretch=3em
\sloppy
\newbox\readingEntry
\hypersetup{pdftitle={Teaching and Source Guide: Agentic AI: Three Operator Cases}}
\begin{document}
{\LARGE\bfseries Teaching and Source Guide}\par
Agentic AI: Three Operator Cases --- introductory lecture. Evidence reviewed 18 October 2026. The full reading guide is separate from the timed main talk. Corporate issuers are listed when no individual author is recorded. Original figures, page locations and image-source credits are included. Attribution alone is not a blanket reuse license.\par
\subsection*{How to teach this follow-on session}
\textbf{Format:} 59 static slides / 120 minutes. Opening 6 minutes; each case 36 minutes including a four-minute exercise; closing 6 minutes.\par
\textbf{Prerequisite:} Use the restored foundations tutorial first. The follow-on starts with concrete operating questions and assumes the terms agent, tool, solver and workflow are already familiar.\par
\textbf{Evidence modes:} Recorded numerical studies, archived genuine LLM replies, and illustrative workflow/version diagrams are labeled separately. No live utility control, new field evaluation, new training or live model call is performed by the deck.\par
\textbf{Preparation:} Present the built Slidev version or the PDF. No model account or installation is needed. Exercises use the supplied slide evidence. Full notebook/source records are available in the source package.\par
\textbf{Facilitation:} For each four-minute exercise: one minute to read, two minutes in pairs, one minute to prepare a decision and one supporting check. The next slide supplies the debrief.\par
\par\medskip
\subsection*{1. Agentic AI: Three Operator Cases}
\textbf{Timing:} 00:00--00:30 (0.5 minutes)\par
\textbf{Slide content:} Voltage support, service restoration and cyber response
120 minutes with three short group exercises
IEEE ComSoc theme / Recorded research examples and operator review\par
\textbf{Teaching notes:} This is the follow-on to the agentic-AI foundations tutorial. We apply the ideas to three engineering decisions, rather than repeating the literature survey. The numerical examples come from the existing GridGuard teaching studies. Their models and assumptions differ: a synthetic distribution feeder, a six-node restoration transport graph, and a four-bus DC estimation fixture. None is a live utility deployment or a permission to operate equipment. All slides belong to the session. The three exercises require discussion and a decision, not software installation.\par
\textbf{[B01] IEEE Communications Society. IEEE Communications Society: official resources and network presentation template. Reviewed 17 October 2026.:} \href{https://www.comsoc.org/about/brand-experience/resources}{\small\ttfamily https:/\allowbreak{}/\allowbreak{}www.\allowbreak{}comsoc.\allowbreak{}org/\allowbreak{}about/\allowbreak{}brand-\allowbreak{}experience/\allowbreak{}resources}\par
\textbf{[W00] GridGuard project teaching materials. Classroom Labs: voltage support, six-node restoration and four-bus cyber response. Existing local GridGuard workshop; recorded numerical examples and source code. 2026; inspected 18 October 2026.:} \href{../research/workshop_evidence.md}{\small\ttfamily research/\allowbreak{}workshop\_\allowbreak{}evidence.\allowbreak{}md}\par
\par\medskip
\subsection*{2. Three cases test three different kinds of evidence}
\textbf{Timing:} 00:30--02:00 (1.5 minutes)\par
\textbf{Slide content:} Voltage: does the proposed control pass the required model scope?
Restoration: who receives service? Cyber: which information can support a decision?\par
\textbf{Teaching notes:} Ask participants to predict which evidence can fail in each case. Voltage has a model-coverage problem: a candidate can pass a reduced model and fail a broader AC audit. Restoration has an explicit objective and topology problem: maximum total load and priority service can legitimately produce different plans. Cyber has an information-trust problem: a small residual can coexist with a wrong estimated state. These are not interchangeable failure modes, so one universal success percentage is inappropriate. Keep the cases separate while reusing the habit of specifying inputs, tools, tests and the human decision.\par
\textbf{Figure scope:} The same advisory process must answer a different engineering question in each case.\par
\textbf{Editable figure:} diagrams/common\_three\_cases.tex\par
\textbf{[W00] GridGuard project teaching materials. Classroom Labs: voltage support, six-node restoration and four-bus cyber response. Existing local GridGuard workshop; recorded numerical examples and source code. 2026; inspected 18 October 2026.:} \href{../research/workshop_evidence.md}{\small\ttfamily research/\allowbreak{}workshop\_\allowbreak{}evidence.\allowbreak{}md}\par
\par\medskip
\subsection*{3. An operator request becomes a checked study}
\textbf{Timing:} 02:00--03:30 (1.5 minutes)\par
\textbf{Slide content:} The agent interprets the request and chooses allowed tools. Solvers calculate the consequences.
Failed or outdated evidence returns to a bounded repair step, or the study stops.\par
\textbf{Teaching notes:} Read the workflow as responsibilities, not a claim that each box is a separate LLM. Awareness fixes the study state. Planning proposes an allowed action. Simulation produces numerical evidence. Verification tests the request, limits and evidence identity. The optimizer or repair step may suggest a new candidate within a fixed budget. A person receives the resulting advisory packet. This structure is a teaching synthesis informed by the existing classroom implementation. The engineering notebooks use a custom Python supervisor. Their shared responsibilities do not establish that they execute in LangGraph; the separate graph lab demonstrates that runtime independently. A numerical solver or fixed script may perform most of the work. The added walkthroughs expose the actual recorded choices, code calls, output values and pending human decisions. The voltage P3 workflow and the standard restoration/cyber workflow are separate dated studies. Each uses its original custom Python supervisor. The present GridGuard app has a separate LangGraph implementation; these archived records do not test that newer runtime.\par
\textbf{Figure scope:} Proposed advisory workflow. Numerical checks control whether a candidate reaches human review.\par
\textbf{Editable figure:} diagrams/common\_operator\_loop.tex\par
\textbf{[W00] GridGuard project teaching materials. Classroom Labs: voltage support, six-node restoration and four-bus cyber response. Existing local GridGuard workshop; recorded numerical examples and source code. 2026; inspected 18 October 2026.:} \href{../research/workshop_evidence.md}{\small\ttfamily research/\allowbreak{}workshop\_\allowbreak{}evidence.\allowbreak{}md}\par
\textbf{[DEF03] Erik S., Barry Zhang. Building effective agents. Anthropic Engineering. 2024-12-19.:} \href{https://www.anthropic.com/engineering/building-effective-agents}{\small\ttfamily https:/\allowbreak{}/\allowbreak{}www.\allowbreak{}anthropic.\allowbreak{}com/\allowbreak{}engineering/\allowbreak{}building-\allowbreak{}effective-\allowbreak{}agents}\par
\par\medskip
\subsection*{4. Every result must stay attached to its inputs}
\textbf{Timing:} 03:30--05:00 (1.5 minutes)\par
\textbf{Slide content:} Keep the case, model scope, data revision, candidate and limits together.
A matching hash detects a changed packet. It does not prove that a sensor or a model is correct.\par
\textbf{Teaching notes:} Use a concrete example: a voltage result at load multiplier 1.00 cannot support a changed request at 1.10 without checking applicability. Likewise, a restoration capacity change or a different measurement-trust set changes the evidence context. The classroom binding routine hashes canonicalized request, candidate and result contents. That establishes local content consistency only. It is not a digital signature, sensor authentication, or a proof of physical feasibility. The operator should be able to see the relevant identity and limits directly in the report. Do not transfer all details into a cryptographic-looking identifier that nobody can inspect.\par
\textbf{Figure scope:} Illustrative evidence identity. These version labels are teaching examples, not an operational record.\par
\textbf{Editable figure:} diagrams/common\_evidence\_binding.tex\par
\textbf{[W00] GridGuard project teaching materials. Classroom Labs: voltage support, six-node restoration and four-bus cyber response. Existing local GridGuard workshop; recorded numerical examples and source code. 2026; inspected 18 October 2026.:} \href{../research/workshop_evidence.md}{\small\ttfamily research/\allowbreak{}workshop\_\allowbreak{}evidence.\allowbreak{}md}\par
\textbf{[OT01] Keith Stouffer, Michael Pease, CheeYee Tang, Timothy Zimmerman, Victoria Pillitteri, Suzanne Lightman, Adam Hahn, Stephanie Saravia, Aslam Sherule, Michael Thompson. Guide to Operational Technology (OT) Security. National Institute of Standards and Technology, NIST SP 800-82 Revision 3. September 2023 (final).:} \href{https://doi.org/10.6028/NIST.SP.800-82r3}{\small\ttfamily https:/\allowbreak{}/\allowbreak{}doi.\allowbreak{}org/\allowbreak{}10.\allowbreak{}6028/\allowbreak{}NIST.\allowbreak{}SP.\allowbreak{}800-\allowbreak{}82r3}\par
\par\medskip
\subsection*{5. The three exercises end with a specific next step}
\textbf{Timing:} 05:00--06:00 (1 minutes)\par
\textbf{Slide content:} Work in pairs: one person proposes the next step, the other checks the evidence.
For each case, choose: continue the study, hold the recommendation, or request missing evidence.
Use one number or one failed check to explain the decision. Each discussion lasts four minutes.\par
\textbf{Teaching notes:} Explain the active-learning format before the first case. Give pairs one minute to read, two minutes to decide and one minute to prepare a short justification. The debrief follows as a separate slide, so the answer is not hidden behind a click or given away on the exercise slide. The choices concern progression through an advisory study, not field dispatch. Encourage participants to name the next responsible tool or human role. A valid hold decision should specify what evidence would allow progress.\par
\textbf{[W00] GridGuard project teaching materials. Classroom Labs: voltage support, six-node restoration and four-bus cyber response. Existing local GridGuard workshop; recorded numerical examples and source code. 2026; inspected 18 October 2026.:} \href{../research/workshop_evidence.md}{\small\ttfamily research/\allowbreak{}workshop\_\allowbreak{}evidence.\allowbreak{}md}\par
\par\medskip
\subsection*{6. Voltage support: return a command with evidence}
\textbf{Timing:} 06:00--07:30 (1.5 minutes)\par
\textbf{Slide content:} Request: use available solar while staying inside voltage, current and inverter limits.
Return an advisory command plus evidence for one synthetic feeder, uhs20\_1247--udt8894.\par
\textbf{Teaching notes:} Start with the operator's work request, not a model brand. The supplied case is one synthetic SMART-DS Greensboro feeder, uhs20\_1247--udt8894. The purpose is to propose fleet inverter settings while checking the declared voltage, current and capability constraints. The command is advisory and does not reach a distribution-management system. The map establishes the study model; it does not establish field accuracy, protection coordination or operational approval. Ask learners to identify the needed inputs: load forecast, previous command, device ratings, available active power assumption, topology/model revision and requested MV/LV coverage. The controller implementation scales inverter ratings, so "available solar" is a declared modeling assumption rather than an irradiance forecast validated here. The twelve-slide path follows request  to  frozen contract  to  proposals  to  numerical validation  to  bounded repair  to  scope audit  to  human review. Evidence: voltage chapter frame V01 and Classroom\_Labs.ipynb cells 2, 4 and 33; source plot figures/feeder.png.\par
\textbf{Figure scope:} Existing SMART-DS synthetic Greensboro feeder; one teaching model, not a live utility circuit.\par
\textbf{Figure asset:} images/feeder.png\par
\textbf{[VCASE] Existing GridGuard workshop / SMART-DS synthetic-feeder extraction. Synthetic Greensboro feeder uhs20\_1247--udt8894: model map and declared study scope. Existing teaching snapshot reviewed 18 October 2026.:} \href{https://data.openei.org/submissions/2981}{\small\ttfamily https:/\allowbreak{}/\allowbreak{}data.\allowbreak{}openei.\allowbreak{}org/\allowbreak{}submissions/\allowbreak{}2981}\par
\textbf{[VGUIDE] GridGuard teaching materials. Instructor guide: a coherent 180-minute teaching session. Existing local workshop snapshot, reviewed 18 October 2026.:} \href{../research/source_VGUIDE.html}{\small\ttfamily research/\allowbreak{}source\_\allowbreak{}VGUIDE.\allowbreak{}html}\par
\par\medskip
\subsection*{7. Freeze what the two control numbers mean}
\textbf{Timing:} 07:30--09:30 (2 minutes)\par
\textbf{Slide content:} The fleet command is (a,b): active participation and reactive injection.
Use declared voltage/current limits; inverter capability alone does not establish feeder feasibility.\par
\textbf{Teaching notes:} Define S\_i as each inverter's kVA rating. The implementation requests p\_i=a S\_i kW and q\_i=b S\_i kvar using one dimensionless fleet command. Negative b means reactive absorption under the injection convention. The action validator requires 0 <= a <= 1, |b| <= 0.44 and a\textasciicircum{}2+b\textasciicircum{}2 <= 1. These are implemented teaching bounds, not universal utility settings. Network checks separately require modeled voltage magnitudes in [0.95,1.05] p.u. and current no greater than the scenario's rating\_fraction times nominal current allowance. "100\% allowed-current use" means the derated allowance is fully used, not necessarily 100\% nameplate current. Freeze the backend, forecast, rating fraction, previous action and model/request revision before interpreting a result. A small negative b may help high voltage in this case, but its direction and feasibility still need a model solve. The illustrated capability region is clipped to the implemented reactive band. The marker uses rounded MPC coordinates and does not prove network feasibility. Evidence: continuous\_control.py \_control, map\_action, ContinuousPlant.solve and verify\_control; Classroom\_Labs.ipynb L1 cell 4.\par
\textbf{Figure scope:} Native teaching diagram of the implemented fleet capability region; the point is the recorded MPC command.\par
\textbf{Editable figure:} diagrams/voltage\_capability.tex\par
\textbf{[VCTRL] GridGuard teaching materials. Continuous fleet P/Q control and classroom validation/readback helpers. Local source reviewed 18 October 2026.:} \href{../research/source_VCTRL.html}{\small\ttfamily research/\allowbreak{}source\_\allowbreak{}VCTRL.\allowbreak{}html}\par
\textbf{[VNOTE] GridGuard teaching materials. Agentic AI for grid operators --- Classroom Labs. Executed 17 October 2026.:} \href{../research/source_VNOTE.html}{\small\ttfamily research/\allowbreak{}source\_\allowbreak{}VNOTE.\allowbreak{}html}\par
\par\medskip
\subsection*{8. Numerical tools propose; the language model selects a trial}
\textbf{Timing:} 09:30--11:30 (2 minutes)\par
\textbf{Slide content:} MPC solves three forecast snapshots and proposes its first move.
SAC proposes from saved trained weights. An LLM selects a permitted trial after numerical checks.\par
\textbf{Teaching notes:} MPC means model predictive control. Here it is a local sensitivity optimizer over H=3 quasi-static power-flow snapshots, linked by a command-movement penalty and discount 0.98. Only its first action enters the response-selection experiment. This is not a certified global optimum, dynamic stability study or robust recursive-feasibility proof. SAC means soft actor-critic: a model-free RL method. The supplied actor was trained previously using rewards from a declared OpenDSS response-surface approximation at nominal current allowance. Classroom inference uses saved weights without retraining; current derating is a transfer condition. Both tools produce candidate numerical commands. A fresh simulation and fail-closed checks decide whether a candidate can become the incumbent. Only then do actual local language models request a bounded refinement from the allowed menu. They do not solve the power flow, relax limits or grant approval. Cite the local implementation for every measured result; the MPC book and SAC paper explain the algorithm families, not performance of this feeder. Evidence: continuous\_control.py optimize\_mpc/sac\_proposal, pilot README Roles and authority and Planning horizon versus first-action checks.\par
\textbf{Figure scope:} Native workflow for the existing implementation; changed commands return to the solver.\par
\textbf{Editable figure:} diagrams/voltage\_proposal\_flow.tex\par
\textbf{[VCTRL] GridGuard teaching materials. Continuous fleet P/Q control and classroom validation/readback helpers. Local source reviewed 18 October 2026.:} \href{../research/source_VCTRL.html}{\small\ttfamily research/\allowbreak{}source\_\allowbreak{}VCTRL.\allowbreak{}html}\par
\textbf{[VMPC] James B. Rawlings; David Q. Mayne; Moritz M. Diehl. Model Predictive Control: Theory, Computation, and Design, second edition. Nob Hill Publishing. 2017.:} \href{https://sites.engineering.ucsb.edu/~jbraw/mpc/MPC-book-2nd-edition-1st-printing.pdf}{\small\ttfamily https:/\allowbreak{}/\allowbreak{}sites.\allowbreak{}engineering.\allowbreak{}ucsb.\allowbreak{}edu/\allowbreak{}\textasciitilde{}jbraw/\allowbreak{}mpc/\allowbreak{}MPC-\allowbreak{}book-\allowbreak{}2nd-\allowbreak{}edition-\allowbreak{}1st-\allowbreak{}printing.\allowbreak{}pdf}\par
\textbf{[VSAC] Tuomas Haarnoja; Aurick Zhou; Pieter Abbeel; Sergey Levine. Soft Actor-Critic: Off-Policy Maximum Entropy Deep Reinforcement Learning with a Stochastic Actor. Proceedings of the 35th International Conference on Machine Learning, PMLR 80, pp. 1861--1870. 2018.:} \href{https://proceedings.mlr.press/v80/haarnoja18b.html}{\small\ttfamily https:/\allowbreak{}/\allowbreak{}proceedings.\allowbreak{}mlr.\allowbreak{}press/\allowbreak{}v80/\allowbreak{}haarnoja18b.\allowbreak{}html}\par
\textbf{[VPILOT] GridGuard teaching materials. Voltage response-selection pilot: how to read the evidence. Protocol frozen 27 September 2026; reviewed 18 October 2026.:} \href{../research/source_VPILOT.html}{\small\ttfamily research/\allowbreak{}source\_\allowbreak{}VPILOT.\allowbreak{}html}\par
\par\medskip
\subsection*{9. Compare two freshly checked proposals}
\textbf{Timing:} 11:30--13:30 (2 minutes)\par
\textbf{Slide content:} Fresh L2: MPC (0.960307,-0.111114); SAC (0.340066,+0.014024).
Both pass the MV checks; current-step J is 0.031104 versus 0.436428.\par
\textbf{Teaching notes:} Use the executed default L2 output in Classroom\_Labs.ipynb cell 7, ID classroom-007, with LOAD\_SCALE=1.0 and RATING\_FRACTION=0.80. Its forecast multipliers derive from the saved base context; LOAD\_SCALE=1.0 does not mean every forecast equals one. MPC returns a=0.960307,b=-0.111114 with vmin=1.023874,vmax=1.038709 and maximum allowed-current utilization 73.619901\%. Saved SAC returns a=0.340066,b=+0.014024 with vmin=1.013243,vmax=1.030000 and utilization 51.679336\%. Both pass the reduced MV contract. Their current-step teaching costs are 0.031104 and 0.436428, respectively. The lower-cost proposal is preferred only after constraints are checked. The cost is dimensionless J=(1-a)\textasciicircum{}2+0.05b\textasciicircum{}2+10 mean[(V-1.01)\textasciicircum{}2]+0.05||u-u\_previous||\textasciicircum{}2. It is not dollars or a field reliability score. A controller cost gap does not measure agent-orchestration benefit. The crop selects only the coherent current panel; the full original plot also has separately sorted voltage curves, whose equal horizontal ranks are not identical buses. L2's timer encloses planning, saved-actor inference and checks, not isolated SAC inference.\par
\textbf{Figure scope:} Replotted stored L2 current checks. 100\% is the declared allowed-current limit.\par
\textbf{Figure asset:} images/voltage\_current\_comparison.png\par
\textbf{[VNOTE] GridGuard teaching materials. Agentic AI for grid operators --- Classroom Labs. Executed 17 October 2026.:} \href{../research/source_VNOTE.html}{\small\ttfamily research/\allowbreak{}source\_\allowbreak{}VNOTE.\allowbreak{}html}\par
\textbf{[VCTRL] GridGuard teaching materials. Continuous fleet P/Q control and classroom validation/readback helpers. Local source reviewed 18 October 2026.:} \href{../research/source_VCTRL.html}{\small\ttfamily research/\allowbreak{}source\_\allowbreak{}VCTRL.\allowbreak{}html}\par
\par\medskip
\subsection*{10. Turn three replies into two physical trial choices}
\textbf{Timing:} 13:30--15:30 (2 minutes)\par
\textbf{Slide content:} Three archived replies collapse to two distinct physical actions.
Two votes choose reactive reduction. Selection buys one numerical trial; it does not authorize operation.\par
\textbf{Teaching notes:} Deliberately switch from fresh L2 to archived original test\_01, whose rating\_fraction is 0.7995473229563319 rather than L2's 0.80. The load forecast is [1.0112062191802\allowbreak{}391,1.0052249857\allowbreak{}79734,1.00069057\allowbreak{}41130825], and previous command [0.2696261713643143,0]. This context has MPC incumbent [0.9603066164910677,-0.11111388575888505]. Llama3.1:8b requests INCREASE\_PV\_0025, giving [0.9628066164910677,-0.11111388575888505]. Llama3.2:3b and qwen2.5:7b-instruct both request RELAX\_Q\_10\_PERCENT, giving [0.9603066164910677,-0.10000249718299654]. Group requests by physical tool and parameters before counting them: three public replies form two candidates. The two-vote candidate gets one simulator trial. Neither a majority nor a bandit score is a probability of electrical safety. The actual recorded reasons include unsupported links from Q reduction to curtailment and a "voltage stability" claim; steady-state first-action checks do not establish stability. Printing these stored replies makes no new LLM call. Evidence: contexts.json record id=test\_01, catalog and case; responses.json context\_id=test\_01; Classroom\_Labs.ipynb cell 10.\par
\textbf{Figure scope:} Actual archived test\_01 choices; native diagram groups equal tool-and-parameter records.\par
\textbf{Editable figure:} diagrams/voltage\_canonical\_trial.tex\par
\textbf{[VARCH] GridGuard teaching materials. Archived immutable voltage contexts and actual local Llama/Qwen response records. Recorded September 2026; reviewed 18 October 2026.:} \href{../research/source_VARCH.html}{\small\ttfamily research/\allowbreak{}source\_\allowbreak{}VARCH.\allowbreak{}html}\par
\textbf{[VNOTE] GridGuard teaching materials. Agentic AI for grid operators --- Classroom Labs. Executed 17 October 2026.:} \href{../research/source_VNOTE.html}{\small\ttfamily research/\allowbreak{}source\_\allowbreak{}VNOTE.\allowbreak{}html}\par
\par\medskip
\subsection*{11. Validate the changed command and bind its evidence}
\textbf{Timing:} 15:30--17:30 (2 minutes)\par
\textbf{Slide content:} Archived Q reduction: J=0.031082, V=1.02442--1.03943 p.u., and 72.93\% allowed-current use.
These are MV-model results. Changed commands or inputs require new bound evidence.\par
\textbf{Teaching notes:} The final\_rechecks.json record context\_id=test\_01, operation\_id=RELAX\_Q\_10\_PERCENT reports control [0.9603066164910677,-0.10000249718299654], score 0.031082078060572192, vmin 1.0244175701973093, vmax 1.039431162620944 and maximum allowed-current utilization 72.92980992911413\%. All recorded voltage, line-current, inverter-capability and dispatch-tracking violation counts are zero. This is a newly checked MV first action in its archived original context. Do not subtract its score from a learner-edited L2 case. The physical verifier reconstructs findings from identified finite voltage/current arrays, freshness/convergence, scenario parameters and actual-versus-requested dispatch. The result still needs a separate request/scope check before readback. Bind the exact request, candidate and result so changed load, command or model revision invalidates old evidence. Hash matching demonstrates local identity/consistency; it is not sensor authentication or a guarantee that evidence is trustworthy. A rejected trial does not replace a valid incumbent. Recompute after every repair; a better public explanation changes no electrical evidence. Evidence: exact final\_rechecks record; classroom\_labs.py trial\_saved\_selection, bind\_evidence and bound\_to; notebook L4 cell 13.\par
\textbf{Figure scope:} Recorded test\_01 Q-reduction result, not a result for an edited classroom scenario.\par
\textbf{Editable figure:} diagrams/voltage\_validation\_gate.tex\par
\textbf{[VRECHECK] GridGuard teaching materials. Distinct final first-action numerical rechecks after bounded response selection. Recorded September 2026.:} \href{../research/source_VRECHECK.html}{\small\ttfamily research/\allowbreak{}source\_\allowbreak{}VRECHECK.\allowbreak{}html}\par
\textbf{[VCTRL] GridGuard teaching materials. Continuous fleet P/Q control and classroom validation/readback helpers. Local source reviewed 18 October 2026.:} \href{../research/source_VCTRL.html}{\small\ttfamily research/\allowbreak{}source\_\allowbreak{}VCTRL.\allowbreak{}html}\par
\textbf{[VNOTE] GridGuard teaching materials. Agentic AI for grid operators --- Classroom Labs. Executed 17 October 2026.:} \href{../research/source_VNOTE.html}{\small\ttfamily research/\allowbreak{}source\_\allowbreak{}VNOTE.\allowbreak{}html}\par
\par\medskip
\subsection*{12. 1. Inputs: the recorded OpenDSS case}
\textbf{Timing:} 17:30--19:30 (2 minutes)\par
\textbf{Slide content:} P3 uses full-feeder OpenDSS, an 80\% current allowance and three forecast loads.
The model chooses allowed steps. Saved SAC and calculated MPC supply candidate controls.\par
\textbf{Teaching notes:} Explicitly start a separate, internally coherent recorded walkthrough. Do not carry over test\_01 controls, costs or response-selection results. The source is workshop\_run\_opendss.json, agentic, case P3\_current\_reserve, recorded 2026-09-27T16:12:29.530248+00:00. Forecast multipliers are [1.00,0.94,1.06], previous command [0.25,0], rating\_fraction=0.80, and modeled voltage bounds [0.95,1.05] p.u. The workflow receives a saved SAC proposal and a previously calculated three-snapshot MPC plan; it independently solves and checks the selected first actions during this recorded workflow. The structured contract labels the data as synthetic scenario/model-generated evidence, with no live SCADA. Both fixed and agentic supervisors have the same numerical proposals and permit at most two candidate evaluations: SAC first and one eligible MPC fallback. The language model's role is to request a permitted next numerical check. The engineering runtime is a custom Python supervisor; it is not the current GridGuard desktop LangGraph implementation. Reading the saved dashboard or notebook now makes no new LLM call or physical command.\par
\textbf{Figure scope:} Actual dashboard Awareness excerpt, P3 OpenDSS, 27 September 2026; separate from test\_01.\par
\textbf{Figure asset:} images/voltage\_input.png\par
\textbf{[TRACE\_VOLT] GridGuard teaching materials. Recorded OpenDSS voltage workflow: P3 case, actual Llama replies, dashboard and Python source. Local executed tutorial and preserved model-call journals. 27 September 2026.:} \href{../research/source_TRACE_VOLT.html}{\small\ttfamily research/\allowbreak{}source\_\allowbreak{}TRACE\_\allowbreak{}VOLT.\allowbreak{}html}\par
\par\medskip
\subsection*{13. 2. Agent: request the MPC fallback}
\textbf{Timing:} 19:30--21:30 (2 minutes)\par
\textbf{Slide content:} Planner chooses VERIFY\_SAC\_FIRST; that proposal then fails the thermal check.
Recovery chooses USE\_MPC\_PLAN. The final explanation asks the operator to REVIEW\_EVIDENCE.\par
\textbf{Teaching notes:} All three replies are genuine records in continuous\_llama\allowbreak{}\_opendss\_2026092\allowbreak{}7T161223523552Z.\allowbreak{}json and are matched to this workflow. Planner, call 1 at 16:12:23.523552 UTC, returns VERIFY\_SAC\_FIRST: 'A numerical SAC proposal is available; its fresh first-action check has not yet been requested.' Its context contains the immutable contract, numerical MPC plan, SAC proposal, current stage, protocol and permitted actions. After numerical thermal rejection, Recovery, call 2 at 16:12:24.665244 UTC, receives failed\_candidate, failed\_checks=['thermal'], replacement\_plan and the same contract. It returns USE\_MPC\_PLAN: 'The MPC plan has been generated and is eligible for a numerical check. This is the next allowed step in the recovery process.' Explanation, call 3 at 16:12:25.646676 UTC, returns REVIEW\_EVIDENCE after the MPC check passes. The client times are 0.9879972, 0.8466333 and 0.9853947 seconds. These replies select bounded transitions; they do not produce the setpoints, prove feasibility, alter the 80\% limit or acknowledge for an operator. The resulting fixed and agentic final controls are the same in this recorded pair. No general performance improvement is demonstrated by adding these calls.\par
\textbf{Figure scope:} Actual dashboard Recovery reply after the thermal rejection; no new model call was made for this slide.\par
\textbf{Figure asset:} images/voltage\_route.png\par
\textbf{[TRACE\_VOLT] GridGuard teaching materials. Recorded OpenDSS voltage workflow: P3 case, actual Llama replies, dashboard and Python source. Local executed tutorial and preserved model-call journals. 27 September 2026.:} \href{../research/source_TRACE_VOLT.html}{\small\ttfamily research/\allowbreak{}source\_\allowbreak{}TRACE\_\allowbreak{}VOLT.\allowbreak{}html}\par
\par\medskip
\subsection*{14. 3. Code: check the chosen command}
\textbf{Timing:} 21:30--23:30 (2 minutes)\par
\textbf{Slide content:} USE\_MPC\_PLAN selects the existing numerical candidate; the agent does not invent new setpoints.
The loop solves that exact command with OpenDSS and calls verify\_control before returning a recommendation.\par
\textbf{Teaching notes:} This excerpt displays the two original statements from src/continuous\_workflow.py lines231--232. Long arguments are wrapped for readability; the builder verifies the same Python syntax tree. plant.solve receives the exact candidate control, first forecast load and current allowance. verify\_control then checks the result. The following original lines233--234 also enforce backend/control identity. Earlier in this same recorded run, Llama selected VERIFY\_SAC\_FIRST, the fresh SAC check failed thermal limits, and Recovery selected USE\_MPC\_PLAN. The bounded loop then checked the existing numerical MPC candidate. The model neither writes setpoints nor changes limits. The statement excerpt is part of a function, not a standalone program. Full original source and the source/call journals are linked under TRACE\_VOLT.\par
\textbf{Figure scope:} Exact source excerpt: the allowed recovery choice returns to the existing simulation-and-verification loop.\par
\textbf{Exact source excerpt:} src/continuous\_workflow.py, lines 231--232. Source statements; line breaks added for readability.
point = plant.solve(
    candidate['control'],
    load\_mult=case.forecast\_loads[0],
    rating\_fraction=case.rating\_fraction,
)
failures = verify\_control(point, case)\par
\textbf{Source alignment:} research/code\_manifest.json records source hashes, original line ranges and syntax-tree equivalence checks. research/code/ contains the readable display fragment. The complete original file is preserved in research/evidence/src/.\par
\textbf{[TRACE\_VOLT] GridGuard teaching materials. Recorded OpenDSS voltage workflow: P3 case, actual Llama replies, dashboard and Python source. Local executed tutorial and preserved model-call journals. 27 September 2026.:} \href{../research/source_TRACE_VOLT.html}{\small\ttfamily research/\allowbreak{}source\_\allowbreak{}TRACE\_\allowbreak{}VOLT.\allowbreak{}html}\par
\par\medskip
\subsection*{15. 4. Result: SAC fails; MPC passes}
\textbf{Timing:} 23:30--25:30 (2 minutes)\par
\textbf{Slide content:} SAC fails at 111.97\% of allowed current. The MPC fallback passes at 99.60\%.
The checked proposal is P = 2387.16 kW and Q = +325.74 kvar; Vmax = 1.03503 p.u.\par
\textbf{Teaching notes:} The saved SAC action is [0.33693869411945343,0.030701177120208742]. OpenDSS converges, voltage is 1.0025570904--1.0375340016 p.u., but current use is 111.9666069740\% of the allowed current, with two line-current findings. It is rejected for thermal limits. The MPC fallback is [0.3040791833492598,0.04149291669866676]. A new selected-first-action solve gives P=2387.1573822298833 kW, Q=325.7379256164953 kvar, Vmin=1.001782946227172 and Vmax=1.03502551840022 p.u., current use=99.60441047658344\% of the allowed current, and zero voltage/current/inverter/dispatch-tracking findings. Nominal-current utilization is 79.68352838126675\%, which becomes 99.6044\% when divided by the 80\% allowance. Keep these denominators explicit. The workflow contains two fresh candidate evaluations and one rejection. Its recorded interval is 3.1104543 seconds, excluding the supplied precomputed controllers/training. The paired fixed workflow reaches the same final MPC control. These results establish what the numerical tools checked in one synthetic full-feeder snapshot, not a benefit caused by LLM routing, dynamic stability or protection adequacy.\par
\textbf{Figure scope:} Recorded P3 first-action results. 100\% denotes the declared current allowance, which is 80\% of the nominal rating.\par
\textbf{Figure asset:} images/voltage\_result.png\par
\textbf{[TRACE\_VOLT] GridGuard teaching materials. Recorded OpenDSS voltage workflow: P3 case, actual Llama replies, dashboard and Python source. Local executed tutorial and preserved model-call journals. 27 September 2026.:} \href{../research/source_TRACE_VOLT.html}{\small\ttfamily research/\allowbreak{}source\_\allowbreak{}TRACE\_\allowbreak{}VOLT.\allowbreak{}html}\par
\par\medskip
\subsection*{16. 5. Operator: review the exact recommendation}
\textbf{Timing:} 25:30--27:30 (2 minutes)\par
\textbf{Slide content:} Read back P = 2387.16 kW, Q = +325.74 kvar, the 80\% current allowance and full MV/LV scope.
Then acknowledge this exact advisory record or request review. The saved human step is still PENDING.\par
\textbf{Teaching notes:} The recorded final status is MODEL\_RECOMMEND. Its last trace event is Operator acknowledgment, PENDING: acknowledged=False, dispatch\_authorized=False and equipment\_operated=False. The recommendation identity is 63408d6a0cdca791\allowbreak{}ff902746f97ef2a2\allowbreak{}d1c27d1f9456846b\allowbreak{}74ccfcacd6f97fb1\allowbreak{}. That binds the displayed case, action, modeled results and contract; it is not authenticated identity or proof that sensors are correct. The original Workshop\_Operator\_Dashboard.html offers a local P/Q readback, requested-current-allowance readback, scope/assumption confirmations and reviewer label. A record may be acknowledged only if its stored numerical evidence is internally consistent and the readback matches. For this result the displayed P/Q is 2387.16 kW and 325.74 kvar, with an 80\% nominal-current allowance. Viewing this screenshot does not acknowledge anything. The screenshot must retain the pending state and blank readback controls. Any later dashboard acknowledgment would be a new unauthenticated local teaching record, not part of the archived experiment or permission to operate equipment. The model still omits protection, switching transients and operational authorization.\par
\textbf{Figure scope:} Original operator controls, still awaiting human readback. No acknowledgment or equipment operation was performed.\par
\textbf{Figure asset:} images/voltage\_review.png\par
\textbf{[TRACE\_VOLT] GridGuard teaching materials. Recorded OpenDSS voltage workflow: P3 case, actual Llama replies, dashboard and Python source. Local executed tutorial and preserved model-call journals. 27 September 2026.:} \href{../research/source_TRACE_VOLT.html}{\small\ttfamily research/\allowbreak{}source\_\allowbreak{}TRACE\_\allowbreak{}VOLT.\allowbreak{}html}\par
\par\medskip
\subsection*{17. Exercise: is this packet ready for the requested scope?}
\textbf{Timing:} 27:30--31:30 (4 minutes)\par
\textbf{Slide content:} The request requires full MV/LV coverage. Choose a disposition, identify missing evidence, and state the next permitted step.\par
\textbf{Teaching notes:} Keep the answer hidden until the following debrief slide. Give pairs four minutes: one person argues for a disposition and the other audits the evidence. The figure contains the exact archived test\_01 Q-reduction result, rounded for display. The full MV/LV request is a hypothetical operator requirement used to examine scope matching, not a claim that the archived experiment dispatched a full-feeder action. Ask learners to distinguish three questions: are the displayed MV numbers within the teaching bounds, does that cover the requested network, and is any equipment action authorized? Answer for the teacher: the MV physical checks pass, but the packet remains HOLD for the requested full MV/LV scope because the evidence covers only MV phases with ideal LV aggregation. A same-action full-feeder OpenDSS audit is needed. If that action fails, use an allowed numerical replan/repair and obtain new checks; do not reuse old evidence or ask the language model to declare approval. Acknowledgment cannot fill a missing scope check. No additional simulation has been run for this exercise. Accept any answer that clearly separates physical pass, scope adequacy and human authority. Evidence: recorded metrics in final\_rechecks.json; operator\_packet scope equality rule; Classroom\_Labs.ipynb cell 33.\par
\textbf{Figure scope:} Four minutes in pairs. Metrics are recorded; the full-scope operator request is the teaching prompt.\par
\textbf{Editable figure:} diagrams/voltage\_exercise\_packet.tex\par
\textbf{[VRECHECK] GridGuard teaching materials. Distinct final first-action numerical rechecks after bounded response selection. Recorded September 2026.:} \href{../research/source_VRECHECK.html}{\small\ttfamily research/\allowbreak{}source\_\allowbreak{}VRECHECK.\allowbreak{}html}\par
\textbf{[VNOTE] GridGuard teaching materials. Agentic AI for grid operators --- Classroom Labs. Executed 17 October 2026.:} \href{../research/source_VNOTE.html}{\small\ttfamily research/\allowbreak{}source\_\allowbreak{}VNOTE.\allowbreak{}html}\par
\par\medskip
\subsection*{18. Debrief: passing one model is not enough}
\textbf{Timing:} 31:30--33:30 (2 minutes)\par
\textbf{Slide content:} The request remains HOLD: MV-only evidence does not meet full MV/LV scope.
Audit the same action in OpenDSS; if it fails, replan and verify a new command.
Operator readback records an advisory packet; it cannot dispatch equipment.\par
\textbf{Teaching notes:} Reveal the answer after pairs have reported. First acknowledge that the displayed MV voltage and current checks pass: this avoids teaching that an unfavorable decision means the calculations were wrong. The mismatch is coverage. A request for full supplied MV/LV needs evidence from that scope. A same-action broader-model audit asks whether this exact command survives changed equations and constraint coverage. A separate AC MPC solve asks what a controller using that broader model would recommend; it changes the action and therefore needs a new final check. Do not call that change an LLM improvement. Keep the failed trial and scope limitation in the final packet. In the actual classroom L12 code, the full-scope packet has evidence\_matches=True and engineering\_pass=True but scope\_matches=False, so it remains HOLD and cannot be acknowledged. For an explicitly MV-scoped teaching request, the same bound result can be READY\_FOR\_READBACK. Wrong recommendation ID or missing operator name rejects readback. A valid local acknowledgment still has equipment\_dispatch=False. This is a semantic guard in a teaching implementation, not a certified live control-room permission system. Evidence: Classroom\_Labs.ipynb cell 33 stored output and src/classroom\_labs.py operator\_packet/acknowledge.\par
\textbf{[VNOTE] GridGuard teaching materials. Agentic AI for grid operators --- Classroom Labs. Executed 17 October 2026.:} \href{../research/source_VNOTE.html}{\small\ttfamily research/\allowbreak{}source\_\allowbreak{}VNOTE.\allowbreak{}html}\par
\textbf{[VCTRL] GridGuard teaching materials. Continuous fleet P/Q control and classroom validation/readback helpers. Local source reviewed 18 October 2026.:} \href{../research/source_VCTRL.html}{\small\ttfamily research/\allowbreak{}source\_\allowbreak{}VCTRL.\allowbreak{}html}\par
\par\medskip
\subsection*{19. Check the same actions against broader model coverage}
\textbf{Timing:} 33:30--35:30 (2 minutes)\par
\textbf{Slide content:} On 12 held-out cases, selected MV actions pass 12/12 MV checks and 0/12 same-action AC audits.
Broader equations and network coverage changed the result; better prose cannot repair missing model scope.\par
\textbf{Teaching notes:} Use the historical twelve-context comparison, not the default L5 notebook cell as evidence for all twelve outcomes. The primary frozen learned selector and majority both have final\_model\_passes=12 and final\_ac\_passes=0 on the MV test split in selector\_analysis.json. The exact retained actions were independently audited in full-feeder OpenDSS. AC audit outcomes were excluded from MV LLM inputs, acceptance and selector reward; do not imply that the MV selector anticipated them. The broad model changes both equations and covered network/constraints, so this is not an isolated numerical-accuracy test. For a concrete initial-action example, test\_01's original MV-MPC command has MV vmax=1.0387087638 and allowed-current utilization=73.6615825\%, yet its same-action AC audit reports vmax=1.0898278087, allowed-current utilization=347.2530224\%, seventeen high-voltage findings and eleven line-current findings. That example is the INITIAL\_MPC audit, not the later Q-reduction record. Independent AC MPC replanning gives first actions passing 12/12 AC checks in these recorded contexts. Neither result is a field safety certificate or a claim of global AC optimality. Evidence: selector\_analysis test summaries and ac\_transfer summaries; initial\_model\_audits test\_01/INITIAL\_MPC; heldout\_model\_comparison.json.\par
\textbf{Figure scope:} Separately recorded 12-case pilot. MV phases with ideal LV aggregation versus the supplied MV/LV model.\par
\textbf{Editable figure:} diagrams/voltage\_model\_scope.tex\par
\textbf{[VANAL] GridGuard teaching materials. Frozen response-selector comparison: held-out MV test and paired full-feeder AC transfer. Protocol frozen 27 September 2026.:} \href{../research/source_VANAL.html}{\small\ttfamily research/\allowbreak{}source\_\allowbreak{}VANAL.\allowbreak{}html}\par
\textbf{[VAC] GridGuard teaching materials. Paired same-action full-feeder audit and independent OpenDSS MPC first-action plans. Recorded September 2026.:} \href{../research/source_VAC.html}{\small\ttfamily research/\allowbreak{}source\_\allowbreak{}VAC.\allowbreak{}html}\par
\textbf{[VPILOT] GridGuard teaching materials. Voltage response-selection pilot: how to read the evidence. Protocol frozen 27 September 2026; reviewed 18 October 2026.:} \href{../research/source_VPILOT.html}{\small\ttfamily research/\allowbreak{}source\_\allowbreak{}VPILOT.\allowbreak{}html}\par
\par\medskip
\subsection*{20. The historical selector ties a simpler voting rule}
\textbf{Timing:} 35:30--37:30 (2 minutes)\par
\textbf{Slide content:} Vote and frozen LinUCB tie on all 12 MV test cases: mean J=0.03398479.
Fixed numerical refinement gives 0.03400656. This pilot does not establish learned-selection superiority.\par
\textbf{Teaching notes:} LinUCB is a linear contextual bandit: it scores which permitted operation to test from context and canonical-action features. It does not train the LLMs or evaluate private reasoning. There are 24 training contexts, 8 validation contexts and 12 held-out tests. The primary selector has 24 selected-feedback updates and remains frozen during validation/test. The actual test summaries give majority\_vote and learned\_selector identical mean final objective 0.0339847924178626. paired\_test\_comparison against majority\_vote reports learned\_better=0, ties=12 and learned\_worse=0. The predeclared fixed numerical refinement is INCREASE\_PV\_0025, with mean 0.0340065598125857. A posthoc fixed Q-reduction rule reproduces the vote/learner choices without LLM calls; label that rule diagnostic, not a predeclared baseline. The tiny cost improvement over one fixed action does not establish a general learned or language-model benefit. All common-pool selectors incur three response calls. Reconstructed serial mean time is 8.119085 seconds for the learned selector versus 3.326839 for the fixed numerical method; these are sums of recorded components, not a measured online control-room transaction. Five training orders reuse the same contexts and do not create five independent datasets. Evidence: selector\_analysis.json protocol, test.summaries and paired\_test\_comparison; README Latency and limits.\par
\textbf{Figure scope:} Historical single-feeder test split; lower dimensionless current-step J is preferred after the MV gate.\par
\textbf{Editable figure:} diagrams/voltage\_selector\_tie.tex\par
\textbf{[VANAL] GridGuard teaching materials. Frozen response-selector comparison: held-out MV test and paired full-feeder AC transfer. Protocol frozen 27 September 2026.:} \href{../research/source_VANAL.html}{\small\ttfamily research/\allowbreak{}source\_\allowbreak{}VANAL.\allowbreak{}html}\par
\textbf{[VARCH] GridGuard teaching materials. Archived immutable voltage contexts and actual local Llama/Qwen response records. Recorded September 2026; reviewed 18 October 2026.:} \href{../research/source_VARCH.html}{\small\ttfamily research/\allowbreak{}source\_\allowbreak{}VARCH.\allowbreak{}html}\par
\par\medskip
\subsection*{21. Replan with the model that covers the request}
\textbf{Timing:} 37:30--39:30 (2 minutes)\par
\textbf{Slide content:} Archived test\_01 AC MPC proposes (0.306052,+0.027697); Vmax=1.03300 p.u. and allowed-current use=99.87\%.
All 12 AC-MPC first actions pass the AC gate; raw SAC passes 1/12.\par
\textbf{Teaching notes:} This single-panel original plot is figures/selector\_E6.png. It compares the active component a of the original MV LinDistFlow incumbent with an independent full-feeder OpenDSS MPC first action for each of the same twelve contexts. It is not a percent-error plot or an LLM-selector comparison. For archived test\_01, heldout\_model\_comparison.json gives AC first action [0.3060522326419636,+0.027696908967886763], vmin=1.0005276166702253, vmax=1.0330024384102277 and maximum allowed-current utilization=99.87273873689041\%; the check passes with no recorded violations. All twelve independently AC-planned first actions pass their snapshot gate. Raw saved SAC passes only one of the twelve full AC checks. The AC transfer reuses already computed MPC plans and marks the planning time as reused; it does not retrain or tune the selector. Its final mean J remains 0.4713848786592687 for both majority and learner, compared with fixed numerical 0.47112137693022893 and uniform response 0.47110513098466494. Thus the learned selector does not beat those simpler comparators here. Do not directly rank MV and AC objective means: the voltage sets and covered constraints differ. A passing first action proves neither full-horizon feasibility nor stability. Evidence: heldout\_model\_comparison array test\_01 opendss\_mpc; selector\_analysis ac\_transfer.summaries mpc\_only/raw\_sac and comparators.\par
\textbf{Figure scope:} Replotted paired command records. Each curve uses a separately planned control; this is not a solver-error plot.\par
\textbf{Figure asset:} images/voltage\_scope\_commands.png\par
\textbf{[VAC] GridGuard teaching materials. Paired same-action full-feeder audit and independent OpenDSS MPC first-action plans. Recorded September 2026.:} \href{../research/source_VAC.html}{\small\ttfamily research/\allowbreak{}source\_\allowbreak{}VAC.\allowbreak{}html}\par
\textbf{[VANAL] GridGuard teaching materials. Frozen response-selector comparison: held-out MV test and paired full-feeder AC transfer. Protocol frozen 27 September 2026.:} \href{../research/source_VANAL.html}{\small\ttfamily research/\allowbreak{}source\_\allowbreak{}VANAL.\allowbreak{}html}\par
\par\medskip
\subsection*{22. Return one scoped packet for human review}
\textbf{Timing:} 39:30--42:00 (2.5 minutes)\par
\textbf{Slide content:} A human acknowledges a specific checked packet.
Scope mismatch, a wrong ID or missing operator name keeps readback blocked.\par
\textbf{Teaching notes:} Ask a learner to narrate the complete operator path: state the request; freeze case and allowed actions; obtain MPC/SAC proposals; calculate fresh evidence; verify every physical bound; test any permitted refinement; audit the requested model scope; and return the exact evidence packet. The packet must preserve the original command, source/model identity, failed trials, current/voltage margins, limitations and a recommendation identifier. A human reviews the specific packet and reads back that identifier with an operator name. In the implementation, valid acknowledgment changes only the local advisory record. It cannot invoke a live DMS, change a setpoint or overrule a failed scope check. The actual L12 output shows a full MV/LV request remaining HOLD despite engineering\_pass=True under MV scope, while an explicitly MV teaching request can be acknowledged. A wrong ID is rejected; a correct readback still leaves equipment\_dispatch=False. Hashes establish consistency of the local request/candidate/result binding, not external authenticity. End with the measured lesson: numerical tools supplied settings; language supplied bounded trial instructions; the MV learner tied voting; a broader model changed admissibility. No general advantage of more LLMs, a learned selector or this teaching architecture has been demonstrated. Evidence: notebook cell 33 and source operator\_packet/acknowledge.\par
\textbf{Figure scope:} Recorded L12 review logic: acknowledgment remains a local simulation record with no equipment dispatch.\par
\textbf{Editable figure:} diagrams/voltage\_operator\_packet.tex\par
\textbf{[VNOTE] GridGuard teaching materials. Agentic AI for grid operators --- Classroom Labs. Executed 17 October 2026.:} \href{../research/source_VNOTE.html}{\small\ttfamily research/\allowbreak{}source\_\allowbreak{}VNOTE.\allowbreak{}html}\par
\textbf{[VCTRL] GridGuard teaching materials. Continuous fleet P/Q control and classroom validation/readback helpers. Local source reviewed 18 October 2026.:} \href{../research/source_VCTRL.html}{\small\ttfamily research/\allowbreak{}source\_\allowbreak{}VCTRL.\allowbreak{}html}\par
\textbf{[VGUIDE] GridGuard teaching materials. Instructor guide: a coherent 180-minute teaching session. Existing local workshop snapshot, reviewed 18 October 2026.:} \href{../research/source_VGUIDE.html}{\small\ttfamily research/\allowbreak{}source\_\allowbreak{}VGUIDE.\allowbreak{}html}\par
\par\medskip
\subsection*{23. Restoration starts with a specific operator request}
\textbf{Timing:} 42:00--44:00 (2 minutes)\par
\textbf{Slide content:} Restore eligible load after isolating fault 1--2.
Use topology T1, the stated tie ratings and critical-load priorities. Return a checked proposal for operator review.\par
\textbf{Teaching notes:} Read the operator request before naming a solver: "Restore eligible load after fault 1--2. Keep the fault isolated. Use the declared ratings and priorities; return a proposal for review." The explicit request is a teaching contract, not an actual dispatch instruction. T1 is a tutorial snapshot label added for this follow-on module. The tools may select a plan, calculate and report evidence; they have no physical switching authority. The six-node transport example is separate from the archived Greensboro feeder and its ampere-based bypass catalog. Spend two minutes identifying the goal, constraints and expected output.\par
\textbf{Figure scope:} Illustrative operator request for the existing six-node classroom fixture.\par
\textbf{Editable figure:} diagrams/restoration\_contract.tex\par
\textbf{[RST01] GridGuard teaching materials. Agentic AI for Grid Operators: restoration teaching chapter, frames R01--R18. Local existing workshop teaching edition. 17 October 2026 edition.:} \href{../research/source_RST01.html}{\small\ttfamily research/\allowbreak{}source\_\allowbreak{}RST01.\allowbreak{}html}\par
\textbf{[RST02] GridGuard teaching materials. Classroom Labs: L6 restoration sweep and L7 changed-contract binding. Local executed Jupyter notebook. 17 October 2026 teaching edition; inspected18 October2026.:} \href{../research/source_RST02.html}{\small\ttfamily research/\allowbreak{}source\_\allowbreak{}RST02.\allowbreak{}html}\par
\textbf{[RST03] GridGuard teaching materials. Six-node restoration graph, MILP and exhaustive enumeration implementation. Local Python source and saved numerical evidence. Existing17 October2026 teaching edition; inspected18 October2026.:} \href{../research/source_RST03.html}{\small\ttfamily research/\allowbreak{}source\_\allowbreak{}RST03.\allowbreak{}html}\par
\par\medskip
\subsection*{24. Lock the fault, topology and load priorities}
\textbf{Timing:} 44:00--46:00 (2 minutes)\par
\textbf{Slide content:} Fault 1--2 remains unavailable throughout this study.
Node 0 supplies the network. Nodes 2 and 5 carry 90 kW of critical demand. Dashed ties are planning options, not issued switching commands.\par
\textbf{Teaching notes:} Demand is [0,40,60,70,50,30] kW at nodes 0--5. Priority weights are [0,1,10,1,1,10]. Healthy source edge 0--1 has 300 kW capacity; faulted edge 1--2 is unavailable even though its nominal capacity is 250 kW. Sections 2--3, 3--4 and 2--5 each have 250 kW capacity. Tie 1--5 is 100 kW; tie 0--4 is the capacity varied in the experiment. This map fixes a transparent classroom topology; it is not an extracted Greensboro switching diagram. Ask learners to identify the isolated edge and the two critical loads. Fault isolation is an assumed input to this study, not a protection diagnosis performed by the agent.\par
\textbf{Figure scope:} Original teaching schematic of the exact synthetic six-node graph; all demands are kW.\par
\textbf{Editable figure:} diagrams/restoration\_fault\_map.tex\par
\textbf{[RST01] GridGuard teaching materials. Agentic AI for Grid Operators: restoration teaching chapter, frames R01--R18. Local existing workshop teaching edition. 17 October 2026 edition.:} \href{../research/source_RST01.html}{\small\ttfamily research/\allowbreak{}source\_\allowbreak{}RST01.\allowbreak{}html}\par
\textbf{[RST02] GridGuard teaching materials. Classroom Labs: L6 restoration sweep and L7 changed-contract binding. Local executed Jupyter notebook. 17 October 2026 teaching edition; inspected18 October2026.:} \href{../research/source_RST02.html}{\small\ttfamily research/\allowbreak{}source\_\allowbreak{}RST02.\allowbreak{}html}\par
\textbf{[RST03] GridGuard teaching materials. Six-node restoration graph, MILP and exhaustive enumeration implementation. Local Python source and saved numerical evidence. Existing17 October2026 teaching edition; inspected18 October2026.:} \href{../research/source_RST03.html}{\small\ttfamily research/\allowbreak{}source\_\allowbreak{}RST03.\allowbreak{}html}\par
\par\medskip
\subsection*{25. Trace the routes before choosing a pickup plan}
\textbf{Timing:} 46:00--48:00 (2 minutes)\par
\textbf{Slide content:} The route 0--1--5--2 avoids the fault and uses the 100 kW tie.
The other tie is only 40 kW in the first case. A 50 kW load block cannot fit through a 40 kW path alone.\par
\textbf{Teaching notes:} Trace the source-to-critical-node path. The 100 kW tie must carry downstream load assigned behind it; node 1's 40 kW is upstream of that bottleneck. The 40 kW tie alone cannot serve the indivisible 50 kW load at node 4. Local load pickup is all-or-nothing. A bus can be energized for transit while its local load remains disconnected; energized and served are different decisions. The optional arrows illustrate static paths only. They neither tell a crew which device to operate first nor establish interlocks, dead-line checks, protection coordination or equipment status.\par
\textbf{Figure scope:} Highlighted route is a feasible path explanation, not a time-ordered switching sequence.\par
\textbf{Editable figure:} diagrams/restoration\_routes.tex\par
\textbf{[RST02] GridGuard teaching materials. Classroom Labs: L6 restoration sweep and L7 changed-contract binding. Local executed Jupyter notebook. 17 October 2026 teaching edition; inspected18 October2026.:} \href{../research/source_RST02.html}{\small\ttfamily research/\allowbreak{}source\_\allowbreak{}RST02.\allowbreak{}html}\par
\textbf{[RST03] GridGuard teaching materials. Six-node restoration graph, MILP and exhaustive enumeration implementation. Local Python source and saved numerical evidence. Existing17 October2026 teaching edition; inspected18 October2026.:} \href{../research/source_RST03.html}{\small\ttfamily research/\allowbreak{}source\_\allowbreak{}RST03.\allowbreak{}html}\par
\par\medskip
\subsection*{26. More restored kW can mean less critical service}
\textbf{Timing:} 48:00--50:00 (2 minutes)\par
\textbf{Slide content:} At the 40 kW tie setting, priority weighting restores 130 kW total, including all 90 critical kW.
Maximizing total kW restores 140 kW, but only 30 critical kW. The objective changes who receives service.\par
\textbf{Teaching notes:} Behind tie 1--5, critical loads 2 and 5 require 60+30=90 kW. Add 40 kW at node 1 upstream: 130 total and 90 critical. The unweighted choice serves loads 3 and 5 behind that tie: 70+30=100 kW, plus the same upstream 40, giving 140 total but 30 critical. Weights 10 for nodes 2 and 5 produce scores 940 versus 410 weighted kW; these are priority scores, not actual power. A finite weight of ten is not universal lexicographic critical-first service. The comparison changes the objective, so it does not measure an LLM benefit. Both methods are conventional exact optimizers on this fixture.\par
\textbf{Figure scope:} Hand-checkable choices on the same graph and constraints; this is a policy comparison.\par
\textbf{Editable figure:} diagrams/restoration\_priority.tex\par
\textbf{[RST01] GridGuard teaching materials. Agentic AI for Grid Operators: restoration teaching chapter, frames R01--R18. Local existing workshop teaching edition. 17 October 2026 edition.:} \href{../research/source_RST01.html}{\small\ttfamily research/\allowbreak{}source\_\allowbreak{}RST01.\allowbreak{}html}\par
\textbf{[RST02] GridGuard teaching materials. Classroom Labs: L6 restoration sweep and L7 changed-contract binding. Local executed Jupyter notebook. 17 October 2026 teaching edition; inspected18 October2026.:} \href{../research/source_RST02.html}{\small\ttfamily research/\allowbreak{}source\_\allowbreak{}RST02.\allowbreak{}html}\par
\textbf{[RST03] GridGuard teaching materials. Six-node restoration graph, MILP and exhaustive enumeration implementation. Local Python source and saved numerical evidence. Existing17 October2026 teaching edition; inspected18 October2026.:} \href{../research/source_RST03.html}{\small\ttfamily research/\allowbreak{}source\_\allowbreak{}RST03.\allowbreak{}html}\par
\par\medskip
\subsection*{27. Ask the solver for a plan tied to this case}
\textbf{Timing:} 50:00--52:00 (2 minutes)\par
\textbf{Slide content:} A mixed-integer linear program (MILP) chooses whole loads and connected branches.
The 40 kW priority plan serves nodes 1, 2 and 5: 130 total, 90 critical kW. Record the chosen edges and per-branch flows.\par
\textbf{Teaching notes:} The highlighted plan uses source\_healthy, tie\_15 and section\_25. Flow is 130 kW on 0--1, 90 on 1--5 and 60 from 5 to 2. The source-file orientation of section\_25 is 2 to 5, so its signed flow is -60 kW. Negative flow is a direction convention, not a violation. Nodes 3 and 4 remain unenergized in this particular minimal-edge plan. The MILP first maximizes priority-weighted pickup, then preserves that optimum while minimizing energized branch count with a small index tie-breaker. This second solve is not a minimum switching-operation sequence. Different feasible or optimal edge sets can exist; the figure states an explicit checked plan rather than claiming uniqueness.\par
\textbf{Figure scope:} One declared feasible priority plan: edges 0--1, 1--5 and 5--2; fault 1--2 stays open.\par
\textbf{Editable figure:} diagrams/restoration\_plan.tex\par
\textbf{[RST02] GridGuard teaching materials. Classroom Labs: L6 restoration sweep and L7 changed-contract binding. Local executed Jupyter notebook. 17 October 2026 teaching edition; inspected18 October2026.:} \href{../research/source_RST02.html}{\small\ttfamily research/\allowbreak{}source\_\allowbreak{}RST02.\allowbreak{}html}\par
\textbf{[RST03] GridGuard teaching materials. Six-node restoration graph, MILP and exhaustive enumeration implementation. Local Python source and saved numerical evidence. Existing17 October2026 teaching edition; inspected18 October2026.:} \href{../research/source_RST03.html}{\small\ttfamily research/\allowbreak{}source\_\allowbreak{}RST03.\allowbreak{}html}\par
\par\medskip
\subsection*{28. Check capacity, source paths and the objective}
\textbf{Timing:} 52:00--54:00 (2 minutes)\par
\textbf{Slide content:} Every picked-up load needs a source-connected path; the energized network must be a tree.
Check branch capacities, active-power balance and fault isolation. Independent enumeration must agree with the small-case MILP objective.\par
\textbf{Teaching notes:} The model uses separate binary variables for energized nodes, picked-up loads and closed edges. Pickup requires an energized node; a selected edge requires energized endpoints. Fictitious commodity flow establishes source connectivity, and the edge count then enforces a tree. Counting edges alone is insufficient. The fictitious flow is not electrical kW. Electrical transport balance and signed capacity inequalities are checked separately. The actual independent routine enumerates pickup and edge choices, verifies source reachability/tree structure and branch demand, and compares the weighted optimum. This is tractable because the fixture is tiny. Agreement validates the finite transport calculation, not the coverage of AC voltage, reactive power, losses, protection or switching dynamics. Wang et al. sections III--IV explain why reachability and edge count are distinct requirements.\par
\textbf{Figure scope:} Numerical verification and evidence binding perform different jobs.\par
\textbf{Editable figure:} diagrams/restoration\_verify.tex\par
\textbf{[RST02] GridGuard teaching materials. Classroom Labs: L6 restoration sweep and L7 changed-contract binding. Local executed Jupyter notebook. 17 October 2026 teaching edition; inspected18 October2026.:} \href{../research/source_RST02.html}{\small\ttfamily research/\allowbreak{}source\_\allowbreak{}RST02.\allowbreak{}html}\par
\textbf{[RST03] GridGuard teaching materials. Six-node restoration graph, MILP and exhaustive enumeration implementation. Local Python source and saved numerical evidence. Existing17 October2026 teaching edition; inspected18 October2026.:} \href{../research/source_RST03.html}{\small\ttfamily research/\allowbreak{}source\_\allowbreak{}RST03.\allowbreak{}html}\par
\textbf{[RSTL01] Ying Wang, Yin Xu, Jiaxu Li, Jinghan He, Xiaojun Wang. On the Radiality Constraints for Distribution System Restoration and Reconfiguration Problems. Author preprint, arXiv:1912.05185v3. 15 March2020(v3); first submitted11 December2019.:} \href{https://arxiv.org/abs/1912.05185v3}{\small\ttfamily https:/\allowbreak{}/\allowbreak{}arxiv.\allowbreak{}org/\allowbreak{}abs/\allowbreak{}1912.\allowbreak{}05185v3}\par
\par\medskip
\subsection*{29. A higher tie rating does not always restore another load}
\textbf{Timing:} 54:00--56:00 (2 minutes)\par
\textbf{Slide content:} At both 80 and 100 kW, the weighted plan serves 200 kW total and all 90 critical kW.
Whole load blocks and available routes create plateaus. These are sampled capacities, not a continuous response curve.\par
\textbf{Teaching notes:} Inspect both panels and the discrete markers. The saved executed L6 cell reports weighted total/critical service 130/90,180/90,200/90,200/90,250/90 at tie capacities 40,60,80,100,120 kW. The unweighted totals/critical values are 140/30,190/30,200/90,200/90,250/90. Fresh exhaustive enumeration of the same source functions reconfirmed those values for this module. The saved notebook's MILP agrees with exhaustive checks; no fresh SciPy MILP is claimed here. At 80/100, one feasible plan supplies nodes 2 and 5 through the 100 kW tie (90 kW), node 3 through tie 0--4 and edge 4--3 (70 kW), and node 1 upstream (40 kW); node 4 is energized for transit but its own 50 kW load is shed. These static paths explain the plateau and are not a switching sequence.\par
\textbf{Figure scope:} Existing computed Classroom Labs L6 figure; both objectives use the same feasible network.\par
\textbf{Figure asset:} images/classroom\_restoration.png\par
\textbf{[RST02] GridGuard teaching materials. Classroom Labs: L6 restoration sweep and L7 changed-contract binding. Local executed Jupyter notebook. 17 October 2026 teaching edition; inspected18 October2026.:} \href{../research/source_RST02.html}{\small\ttfamily research/\allowbreak{}source\_\allowbreak{}RST02.\allowbreak{}html}\par
\textbf{[RST03] GridGuard teaching materials. Six-node restoration graph, MILP and exhaustive enumeration implementation. Local Python source and saved numerical evidence. Existing17 October2026 teaching edition; inspected18 October2026.:} \href{../research/source_RST03.html}{\small\ttfamily research/\allowbreak{}source\_\allowbreak{}RST03.\allowbreak{}html}\par
\textbf{[RST05] GridGuard teaching materials. Classroom restoration capacity sweep: total and critical service. Original computed local scientific figure. 17 October2026 teaching edition; inspected18 October2026.:} \href{../research/source_RST05.html}{\small\ttfamily research/\allowbreak{}source\_\allowbreak{}RST05.\allowbreak{}html}\par
\par\medskip
\subsection*{30. 1. Inputs: the six-node restoration task}
\textbf{Timing:} 56:00--58:00 (2 minutes)\par
\textbf{Slide content:} This recorded six-node run fixes tie 0--4 at 100 kW; the earlier example used 40 kW.
The structured input already names RESTORATION\_MILP as the allowed tool.\par
\textbf{Teaching notes:} Spend the first minute identifying what is supplied rather than inferred. The original Awareness response contains use\_case=restoration, authorized\_tool=RESTORATION\_MILP, the six-node active-power transport scope, backend=separate\_reduced\_teaching\_model, and approved transparent simulation fixtures. It does not contain live utility measurements or a free-text operator request. The tie rating is fixed by numerical\_tool at solve\_restoration\_milp(100.), not selected by Llama. This is the same graph family as our classroom example, with tie 0--4 now at 100 kW rather than the initial 40 kW. Tie 1--5 remains 100 kW; fault 1--2 remains unavailable. Full evidence: C:/Users/f00863x\allowbreak{}/OPRO/output/age\allowbreak{}ntic\_grid\_workfl\allowbreak{}ow\_tutorial/resu\allowbreak{}lts/standard\_age\allowbreak{}nt\_workflow.json\allowbreak{}, runs[5], case=restoration, supervisor=agentic, trace[0]. Recorded 21 September 2026. Matching executed display: Standard\_Formulations.ipynb, zero-based cell 19, id e3659785, execution\_count 8. Actual model journal: results/standard\allowbreak{}\_stage\_20260921T\allowbreak{}143826Z\_42d19a61\allowbreak{}/llama\_calls.jso\allowbreak{}n, call 5 context. The screenshot is a crop of this saved notebook output, not a fresh execution or a utility control-room deployment.\par
\textbf{Figure scope:} Actual Awareness response containing the structured Planning input, 21 September 2026; tie0--4 = 100 kW.\par
\textbf{Figure asset:} images/rest\_input.png\par
\textbf{[TRACE\_REST] GridGuard teaching materials. Recorded six-node restoration workflow: actual Llama routing, MILP evidence and pending human review. Local executed tutorial and preserved model-call journals. 21 September 2026.:} \href{../research/source_TRACE_REST.html}{\small\ttfamily research/\allowbreak{}source\_\allowbreak{}TRACE\_\allowbreak{}REST.\allowbreak{}html}\par
\par\medskip
\subsection*{31. 2. Agent: choose the authorized MILP}
\textbf{Timing:} 58:00--60:00 (2 minutes)\par
\textbf{Slide content:} Planning call 5 returns RESTORATION\_MILP using the local model llama3.1:8b.
Python checks the allowed choice and calls its fixed solver routine. Llama does not calculate the plan.\par
\textbf{Teaching notes:} Read the exact choice before the explanation. Journal call 5 has role Planning and model llama3.1:8b. Its supplied schema allows VOLTVAR\_OPF, RESTORATION\_MILP, WLS\_TRUSTED\_REDISPATCH or REVIEW. The returned choice is RESTORATION\_MILP; the returned reason says the authorized tool is suitable for load pickup and energized-tree selection. Because the authorized tool is already included in the input, this demonstrates bounded routing, not independent tool discovery. The actual provider response is schema-constrained JSON in an Ollama chat message. It is not a native provider tool\_calls message. Python compares the choice with TOOLS[case]; a different family would stop this route before numerical execution. The wrapper, not the language model, fixes the tie capacity and calls the optimizer. Exact evidence: C:/Users/f00863x\allowbreak{}/OPRO/output/age\allowbreak{}ntic\_grid\_workfl\allowbreak{}ow\_tutorial/resu\allowbreak{}lts/standard\_age\allowbreak{}nt\_workflow.json\allowbreak{}, runs[5].trace[1].response.routing; original journal results/standard\allowbreak{}\_stage\_20260921T\allowbreak{}143826Z\_42d19a61\allowbreak{}/llama\_calls.jso\allowbreak{}n, call\_id 5, zero-based index 4. Matching notebook: Standard\_Formulations.ipynb, cell 19, Planning --- SOLVED details. Current src/standard\_agent\_workflow.py has the same SHA-256 recorded in the run metadata.\par
\textbf{Figure scope:} Crop of the original Planning response: actual Llama choice, returned reason, call ID and provenance.\par
\textbf{Figure asset:} images/rest\_route.png\par
\textbf{[TRACE\_REST] GridGuard teaching materials. Recorded six-node restoration workflow: actual Llama routing, MILP evidence and pending human review. Local executed tutorial and preserved model-call journals. 21 September 2026.:} \href{../research/source_TRACE_REST.html}{\small\ttfamily research/\allowbreak{}source\_\allowbreak{}TRACE\_\allowbreak{}REST.\allowbreak{}html}\par
\par\medskip
\subsection*{32. 3. Code: solve and check the pickup plan}
\textbf{Timing:} 60:00--62:00 (2 minutes)\par
\textbf{Slide content:} The actual call is solve\_restoration\_milp(100.): the argument sets tie 0--4 to 100 kW.
The numerical gate requires solver success and agreement with independent enumeration.\par
\textbf{Teaching notes:} The display selects the original statements at src/standard\_agent\_workflow.py lines80 and83. The intervening summary extraction is omitted and disclosed by the noncontiguous line citation. Only line breaks and indentation change; the builder verifies equivalent Python syntax trees. solve\_restoration\_milp(100.) fixes tie0--4 capacity at100 kW in the same six-node graph taught here. accepted requires both MILP success and independent enumeration agreement. This is a graph/transport check, with no AC voltage, protection or temporal switching certificate. Llama only routes the structured input to its already authorized tool. Both fixed and agentic supervisors run the same numerical calculation. Full original function and result/journal are linked under TRACE\_REST.\par
\textbf{Figure scope:} Seven unchanged source lines: fixed numerical call, reported outputs and the limited numerical acceptance test.\par
\textbf{Exact source excerpt:} src/standard\_agent\_workflow.py, lines 80, 83. Source statements; line breaks added for readability.
result = solve\_restoration\_milp(100.)
accepted = bool(
    result['milp']['success']
    and result['agreement']
)\par
\textbf{Source alignment:} research/code\_manifest.json records source hashes, original line ranges and syntax-tree equivalence checks. research/code/ contains the readable display fragment. The complete original file is preserved in research/evidence/src/.\par
\textbf{[TRACE\_REST] GridGuard teaching materials. Recorded six-node restoration workflow: actual Llama routing, MILP evidence and pending human review. Local executed tutorial and preserved model-call journals. 21 September 2026.:} \href{../research/source_TRACE_REST.html}{\small\ttfamily research/\allowbreak{}source\_\allowbreak{}TRACE\_\allowbreak{}REST.\allowbreak{}html}\par
\par\medskip
\subsection*{33. 4. Result: 200 kW total, 90 kW critical}
\textbf{Timing:} 62:00--64:00 (2 minutes)\par
\textbf{Slide content:} The solver serves nodes 1, 2, 3 and 5: 200 kW total, including all 90 critical kW.
Solver success and enumeration agree. The fixed workflow returns the same numerical result without Llama.\par
\textbf{Teaching notes:} Connect the pickup vector [0,1,1,1,0,1] to nodes 1,2,3,5 and their demands 40+60+70+30=200 kW. Critical nodes 2 and 5 contribute 60+30=90 kW. The weighted objective is 1010 weighted kW, a priority score rather than physical power. Selected edges are source\_healthy, section\_34, section\_25, tie\_04 and tie\_15; the reported MILP gap is 0. Solver success and independent\_enumeration\_agreement are both true. Node 4 can be energized for transit while its own 50 kW load is not picked up. The fixed comparator uses identical numerical parameters and returns the same summary: no numerical advantage from Llama is established. Exact evidence: C:/Users/f00863x\allowbreak{}/OPRO/output/age\allowbreak{}ntic\_grid\_workfl\allowbreak{}ow\_tutorial/resu\allowbreak{}lts/standard\_age\allowbreak{}nt\_workflow.json\allowbreak{}, runs[5].summary and trace[4] for the agentic case; runs[4] for the fixed comparator. Matching executed output: Standard\_Formulations.ipynb, cell 19, Planning/Execution proposal and Verification details. Recovery is NOT\_NEEDED with additional\_attempts=0; do not claim agent repair. The evidence covers discrete pickup, source connectivity, radial graph flow and stated capacities. It does not verify AC voltage, protection, field isolation or a switching sequence.\par
\textbf{Figure scope:} Original numerical-plan excerpt. Separate verification records solver success and enumeration agreement.\par
\textbf{Figure asset:} images/rest\_result.png\par
\textbf{[TRACE\_REST] GridGuard teaching materials. Recorded six-node restoration workflow: actual Llama routing, MILP evidence and pending human review. Local executed tutorial and preserved model-call journals. 21 September 2026.:} \href{../research/source_TRACE_REST.html}{\small\ttfamily research/\allowbreak{}source\_\allowbreak{}TRACE\_\allowbreak{}REST.\allowbreak{}html}\par
\par\medskip
\subsection*{34. 5. Operator: the record requires review}
\textbf{Timing:} 64:00--66:00 (2 minutes)\par
\textbf{Slide content:} Llama's Explanation call 6 returns REVIEW despite passing numerical checks. The recorded human acknowledgment remains PENDING.
Proposed next step: inspect the scope, request missing engineering checks and review again before any separate operating authorization.\par
\textbf{Teaching notes:} Read the enum and the returned prose separately. Actual journal call 6 returns choice=REVIEW. Its reason describes successful modeled checks, 200 kW served and 90 kW critical service, then says the results should be reviewed. The wrapper deliberately treats this choice as a hold on its recommendation: src/standard\_agent\_workflow.py lines 134--141 sets accepted=False when report.choice is REVIEW. Final status is REVIEW. This does not prove the numerical plan is infeasible. The same fixed numerical route ends MODEL\_RECOMMEND, which also requires human review and is not dispatch. Exact record: C:/Users/f00863x\allowbreak{}/OPRO/output/age\allowbreak{}ntic\_grid\_workfl\allowbreak{}ow\_tutorial/resu\allowbreak{}lts/standard\_age\allowbreak{}nt\_workflow.json\allowbreak{}, runs[5].trace[6], trace[7] and status. Actual reply: results/standard\allowbreak{}\_stage\_20260921T\allowbreak{}143826Z\_42d19a61\allowbreak{}/llama\_calls.jso\allowbreak{}n, call\_id 6, zero-based index 5. The final recorded human stage says PENDING, acknowledged=false, dispatch\_authorized=false and equipment\_operated=false. No actual human acceptance or rejection is recorded. The second bullet is a proposed operator response for discussion, not a historical human decision. Request whatever current AC, equipment-status, protection and procedural evidence the intended operation requires; approval of a classroom calculation grants no switching authority. This recorded case demonstrates integration and an explanation-interface limitation, not numerical improvement, successful repair or measured operator-time savings.\par
\textbf{Figure scope:} Original Explanation output. The separate acknowledgment record remains PENDING, with no dispatch or equipment operation.\par
\textbf{Figure asset:} images/rest\_review.png\par
\textbf{[TRACE\_REST] GridGuard teaching materials. Recorded six-node restoration workflow: actual Llama routing, MILP evidence and pending human review. Local executed tutorial and preserved model-call journals. 21 September 2026.:} \href{../research/source_TRACE_REST.html}{\small\ttfamily research/\allowbreak{}source\_\allowbreak{}TRACE\_\allowbreak{}REST.\allowbreak{}html}\par
\par\medskip
\subsection*{35. A changed topology or rating makes old evidence stale}
\textbf{Timing:} 66:00--68:00 (2 minutes)\par
\textbf{Slide content:} A changed rating or topology revision invalidates the old evidence binding.
L7 returns True for the original 40 kW request and False after it changes to 20 kW. HOLD and recheck; binding failure alone does not prove infeasibility.\par
\textbf{Teaching notes:} Actual L7 binds request, candidate and result using SHA-256 digests, then changes capacity from 40 to 20 kW and revision from 1 to 2; bound\_to returns True then False. This module additionally executes the same unmodified binding functions on a request containing the full graph snapshot and topology\_revision. It removes unused section\_23 and changes T1/revision1 to T2/revision2; old evidence returns False. The previous plan did not use that branch, illustrating why identity failure does not itself establish physical infeasibility. The explicit graph snapshot is essential: a hash can only bind fields that are included. Hashes do not authenticate the data source, prove measurement freshness or grant operational authority. HOLD the old recommendation and obtain fresh checks under the changed request.\par
\textbf{Figure scope:} L7 changes the rating/revision; the explicit topology-change test is a separately labeled teaching extension.\par
\textbf{Editable figure:} diagrams/restoration\_stale.tex\par
\textbf{[RST02] GridGuard teaching materials. Classroom Labs: L6 restoration sweep and L7 changed-contract binding. Local executed Jupyter notebook. 17 October 2026 teaching edition; inspected18 October2026.:} \href{../research/source_RST02.html}{\small\ttfamily research/\allowbreak{}source\_\allowbreak{}RST02.\allowbreak{}html}\par
\textbf{[RST04] GridGuard teaching materials. Evidence digest, binding and contract matching functions. Local classroom implementation. Existing17 October2026 teaching edition; inspected18 October2026.:} \href{../research/source_RST04.html}{\small\ttfamily research/\allowbreak{}source\_\allowbreak{}RST04.\allowbreak{}html}\par
\textbf{[RST06] GridGuard teaching materials. Fresh restoration enumeration and explicit topology-binding teaching check. Local reproducible teaching evidence. 18 October2026.:} \href{../research/source_RST06.html}{\small\ttfamily research/\allowbreak{}source\_\allowbreak{}RST06.\allowbreak{}html}\par
\par\medskip
\subsection*{36. Give the operator a packet that can be checked and read back}
\textbf{Timing:} 68:00--70:00 (2 minutes)\par
\textbf{Slide content:} Show picked-up loads, selected paths and total/critical service.
Attach topology and rating revisions, candidate identity, solver output and independent checks. Ask the operator to confirm conditions before accepting, revising or rejecting the proposal.\par
\textbf{Teaching notes:} An example readback is: "For topology T1, fault 1--2 isolated, tie 1--5 at 100 kW and tie 0--4 at 40 kW, I understand this proposes serving nodes 1,2,5:130 kW total and 90 critical. Transport and enumeration checks match this candidate. AC/security and switching authorization are outside this packet." The packet also records shed loads, limitations and any unsuccessful checks. A readback is a tutorial communication step, not proof that a field procedure has been completed. Acceptance of the calculation must remain separate from switching authorization. Missing status, failed binding, failed checks or conflicting instructions produce HOLD rather than an action request.\par
\textbf{Figure scope:} Illustrative review packet, not a switching order or a deployed control-room interface.\par
\textbf{Editable figure:} diagrams/restoration\_packet.tex\par
\textbf{[RST01] GridGuard teaching materials. Agentic AI for Grid Operators: restoration teaching chapter, frames R01--R18. Local existing workshop teaching edition. 17 October 2026 edition.:} \href{../research/source_RST01.html}{\small\ttfamily research/\allowbreak{}source\_\allowbreak{}RST01.\allowbreak{}html}\par
\textbf{[RST02] GridGuard teaching materials. Classroom Labs: L6 restoration sweep and L7 changed-contract binding. Local executed Jupyter notebook. 17 October 2026 teaching edition; inspected18 October2026.:} \href{../research/source_RST02.html}{\small\ttfamily research/\allowbreak{}source\_\allowbreak{}RST02.\allowbreak{}html}\par
\textbf{[RST04] GridGuard teaching materials. Evidence digest, binding and contract matching functions. Local classroom implementation. Existing17 October2026 teaching edition; inspected18 October2026.:} \href{../research/source_RST04.html}{\small\ttfamily research/\allowbreak{}source\_\allowbreak{}RST04.\allowbreak{}html}\par
\par\medskip
\subsection*{37. Exercise: choose the service policy and handle a changed case}
\textbf{Timing:} 70:00--74:00 (4 minutes)\par
\textbf{Slide content:} With tie 0--4 at 40 kW, compare loads 2+5 with loads 3+5 behind the 100 kW tie.
Add node 1's 40 kW. Report total and critical service for each choice. The topology changes after verification. Should the old packet be accepted?\par
\textbf{Teaching notes:} Allow the complete four minutes. First minute: identify demands and which loads are critical. Second minute: calculate both service outcomes and explain the preferred policy. Third minute: check a feasible source path and the 100 kW bottleneck. Fourth minute: decide what a topology change does to the evidence. Answers: 2+5 gives90 kW through the tie, plus upstream40 =130 total/90 critical. 3+5 gives100 through the tie, plus40 =140 total/30 critical. The declared priority-weighted policy chooses the first outcome. Fault1--2 cannot be reclosed. A topology change requires HOLD of the old bound packet and a new check; it does not by itself prove the old candidate infeasible. No participant should invent a field switching order from the static map. Keep answers off this exercise slide.\par
\textbf{Figure scope:} Four-minute pair exercise. Work from the map; answers appear in the debrief.\par
\textbf{Editable figure:} diagrams/restoration\_fault\_map.tex\par
\textbf{[RST02] GridGuard teaching materials. Classroom Labs: L6 restoration sweep and L7 changed-contract binding. Local executed Jupyter notebook. 17 October 2026 teaching edition; inspected18 October2026.:} \href{../research/source_RST02.html}{\small\ttfamily research/\allowbreak{}source\_\allowbreak{}RST02.\allowbreak{}html}\par
\textbf{[RST03] GridGuard teaching materials. Six-node restoration graph, MILP and exhaustive enumeration implementation. Local Python source and saved numerical evidence. Existing17 October2026 teaching edition; inspected18 October2026.:} \href{../research/source_RST03.html}{\small\ttfamily research/\allowbreak{}source\_\allowbreak{}RST03.\allowbreak{}html}\par
\textbf{[RST04] GridGuard teaching materials. Evidence digest, binding and contract matching functions. Local classroom implementation. Existing17 October2026 teaching edition; inspected18 October2026.:} \href{../research/source_RST04.html}{\small\ttfamily research/\allowbreak{}source\_\allowbreak{}RST04.\allowbreak{}html}\par
\par\medskip
\subsection*{38. Debrief: keep the priority and the evidence current}
\textbf{Timing:} 74:00--76:00 (2 minutes)\par
\textbf{Slide content:} The priority choice serves 130/90 kW total/critical; the maximum-total choice serves 140/30.
Both can satisfy this transport model while serving different policies. After a topology change, HOLD the old packet and bind fresh evidence.\par
\textbf{Teaching notes:} Spend the first minute asking one pair to explain the arithmetic and the 100 kW tie budget. Use the second minute to explain the stale-evidence response. The difference between130 and140 kW is not a solver error: objectives differ. The original critical weighting yields940 versus410 weighted kW in these example plans. It is not a measured benefit of an LLM. A correct numerical answer from the old request is not current evidence after the request changes. If a learner says "the old plan is definitely infeasible," correct that: changed identity requires rechecking; only fresh engineering evaluation answers feasibility. A finite priority weight does not guarantee universal strict priority.\par
\textbf{Figure scope:} Two-minute debrief: the objective and evidence validity are separate checks.\par
\textbf{Editable figure:} diagrams/restoration\_priority.tex\par
\textbf{[RST02] GridGuard teaching materials. Classroom Labs: L6 restoration sweep and L7 changed-contract binding. Local executed Jupyter notebook. 17 October 2026 teaching edition; inspected18 October2026.:} \href{../research/source_RST02.html}{\small\ttfamily research/\allowbreak{}source\_\allowbreak{}RST02.\allowbreak{}html}\par
\textbf{[RST03] GridGuard teaching materials. Six-node restoration graph, MILP and exhaustive enumeration implementation. Local Python source and saved numerical evidence. Existing17 October2026 teaching edition; inspected18 October2026.:} \href{../research/source_RST03.html}{\small\ttfamily research/\allowbreak{}source\_\allowbreak{}RST03.\allowbreak{}html}\par
\textbf{[RST04] GridGuard teaching materials. Evidence digest, binding and contract matching functions. Local classroom implementation. Existing17 October2026 teaching edition; inspected18 October2026.:} \href{../research/source_RST04.html}{\small\ttfamily research/\allowbreak{}source\_\allowbreak{}RST04.\allowbreak{}html}\par
\textbf{[RST06] GridGuard teaching materials. Fresh restoration enumeration and explicit topology-binding teaching check. Local reproducible teaching evidence. 18 October2026.:} \href{../research/source_RST06.html}{\small\ttfamily research/\allowbreak{}source\_\allowbreak{}RST06.\allowbreak{}html}\par
\par\medskip
\subsection*{39. Stop at the boundary of the model and the operator's authority}
\textbf{Timing:} 76:00--78:00 (2 minutes)\par
\textbf{Slide content:} This fixture checks lossless active-power transport, source connectivity and stated branch capacities.
It does not certify AC voltage, reactive power, losses, protection or a switching sequence.
A reviewed proposal still needs the utility's separate authorization and operating procedure.\par
\textbf{Teaching notes:} Conclude with a scope readback. The agent can preserve the requested priorities, select allowed tools, assemble evidence, detect a stale binding and ask for review. The conventional optimizer chooses a transport plan; no field improvement is established by this case. A larger real restoration model may need three-phase AC checks, cold-load pickup, DER availability, dynamic limits, protection coordination and a time-ordered procedure. Chen et al. is a background reference on sequential restoration, not the source of these six-node numbers. Do not import the separate archived feeder's ampere ratings, served-kW figures or Llama latencies into this module. The next case can ask what happens when the measurements feeding a correct optimizer are wrong.\par
\textbf{[RST01] GridGuard teaching materials. Agentic AI for Grid Operators: restoration teaching chapter, frames R01--R18. Local existing workshop teaching edition. 17 October 2026 edition.:} \href{../research/source_RST01.html}{\small\ttfamily research/\allowbreak{}source\_\allowbreak{}RST01.\allowbreak{}html}\par
\textbf{[RST03] GridGuard teaching materials. Six-node restoration graph, MILP and exhaustive enumeration implementation. Local Python source and saved numerical evidence. Existing17 October2026 teaching edition; inspected18 October2026.:} \href{../research/source_RST03.html}{\small\ttfamily research/\allowbreak{}source\_\allowbreak{}RST03.\allowbreak{}html}\par
\textbf{[RST04] GridGuard teaching materials. Evidence digest, binding and contract matching functions. Local classroom implementation. Existing17 October2026 teaching edition; inspected18 October2026.:} \href{../research/source_RST04.html}{\small\ttfamily research/\allowbreak{}source\_\allowbreak{}RST04.\allowbreak{}html}\par
\textbf{[RSTL02] Bo Chen, Chen Chen, Jianhui Wang, Karen L. Butler-Purry. Sequential Service Restoration for Unbalanced Distribution Systems and Microgrids. IEEE Transactions on Power Systems33(2),1507--1520. March2018.:} \href{https://doi.org/10.1109/TPWRS.2017.2720122}{\small\ttfamily https:/\allowbreak{}/\allowbreak{}doi.\allowbreak{}org/\allowbreak{}10.\allowbreak{}1109/\allowbreak{}TPWRS.\allowbreak{}2017.\allowbreak{}2720122}\par
\par\medskip
\subsection*{40. Freeze the assessment before changing control}
\textbf{Timing:} 78:00--79:30 (1.5 minutes)\par
\textbf{Slide content:} An unusual reading starts an investigation. It does not identify an attacker.
Keep the readings, trust assumptions, topology and model in the same assessment packet.\par
\textbf{Teaching notes:} Set the operator problem: a network estimate suggests that no corrective action is needed, but the provenance of some measurements is unresolved. Before asking for a new control, record the reading set, its time, trust flags, topology, model H, noise scales, operating limits and proposed action. A repeated calculation on changed inputs is a new assessment. The three horizontal tracks make this shared context visible; they do not claim that a hash, an LLM or the workshop code authenticates sensors. In the laboratory, trusted\_response returns HOLD and dispatch=false when an independent trusted state is unavailable. A residual alarm may reflect a bad reading, topology or model error, or an adversarial change. It alone identifies none of these causes. The upcoming Hc example is a defensive negative control on synthetic data, not a live-system test.\par
\textbf{Figure scope:} Proposed defensive assessment discipline for the four-bus teaching fixture.\par
\textbf{Editable figure:} diagrams/cyber\_freeze.tex\par
\textbf{[CY04] GridGuard teaching materials. Standard formulation companion: four-bus DC weighted least squares, residual tests and constrained redispatch. Local reproducible teaching fixture. 21 September 2026.:} \href{../research/source_CY04.html}{\small\ttfamily research/\allowbreak{}source\_\allowbreak{}CY04.\allowbreak{}html}\par
\textbf{[CY07] GridGuard teaching materials. Workshop Instructor Guide and Cyber Teaching Chapter: evidence scope and operator interpretation. Local instructor material. Reviewed 18 October 2026.:} \href{../research/source_CY07.html}{\small\ttfamily research/\allowbreak{}source\_\allowbreak{}CY07.\allowbreak{}html}\par
\par\medskip
\subsection*{41. Eight readings describe three unknown angles}
\textbf{Timing:} 79:30--81:30 (2 minutes)\par
\textbf{Slide content:} The state contains three voltage angles. Eight readings provide five extra consistency checks.
This is a DC model: it does not test voltage magnitudes or reactive power.\par
\textbf{Teaching notes:} Bus 0 is the reference, with angle zero. The unknown vector is x=(theta1,theta2,theta3). The four line susceptance coefficients are 10, 8, 10 and 5 for reference branches 01, 12, 23 and 03. The linear measurement matrix stacks four flow rows F, three injection rows B, and an angle-1 row. Its rank is three. Thus m=8 measurements minus rank(H)=3 gives five residual degrees of freedom. The arrows define algebraic reference directions; actual flow can be negative. In this fixed DC approximation, active flow is linear in angle differences. It omits voltage magnitude, reactive power, losses, protection behavior and AC constraints. Do not transfer this result to the distribution-feeder model used in the voltage case. Ask participants which extra measurements provide consistency information beyond the minimum needed to solve for the angles.\par
\textbf{Figure scope:} Exact topology and measurement count in the supplied four-bus fixture. Arrows define positive flow directions.\par
\textbf{Editable figure:} diagrams/cyber\_network.tex\par
\textbf{[CY04] GridGuard teaching materials. Standard formulation companion: four-bus DC weighted least squares, residual tests and constrained redispatch. Local reproducible teaching fixture. 21 September 2026.:} \href{../research/source_CY04.html}{\small\ttfamily research/\allowbreak{}source\_\allowbreak{}CY04.\allowbreak{}html}\par
\textbf{[CY01] Fred C. Schweppe, J. Wildes. Power System Static-State Estimation, Part I: Exact Model. IEEE Transactions on Power Apparatus and Systems, PAS-89(1), pp. 120--125. January 1970.:} \href{https://doi.org/10.1109/TPAS.1970.292678}{\small\ttfamily https:/\allowbreak{}/\allowbreak{}doi.\allowbreak{}org/\allowbreak{}10.\allowbreak{}1109/\allowbreak{}TPAS.\allowbreak{}1970.\allowbreak{}292678}\par
\par\medskip
\subsection*{42. WLS finds the best common explanation}
\textbf{Timing:} 81:30--84:00 (2.5 minutes)\par
\textbf{Slide content:} A residual is a reading minus the value predicted by the estimated state.
WLS gives more weight to a reading with a smaller stated noise scale.\par
\textbf{Teaching notes:} Start with the scalar picture. Two equally precise readings of one quantity are 9 and 11. Choosing 10 leaves residuals -1 and +1; moving the estimate toward either reading increases the total squared mismatch. WLS extends this idea to the network predictions Hx. It minimizes the sum over readings of ((z\_j-(Hx)\_j)/sigma\_j)\textasciicircum{}2. Here z is the reading vector, H maps angles to expected readings, and sigma\_j is the assumed standard deviation of the reading noise. A smaller sigma creates a larger penalty for the same raw mismatch. It expresses the estimator's stated precision model, not a cybersecurity trust judgment. The supplied code solves the full-rank weighted problem numerically using least squares. Flow noise scales are 0.004 p.u., injection noise scales 0.006 p.u., and the angle-1 noise scale is 0.0005 rad. Do not interpret unequal weights as proof that a channel is authentic.\par
\textbf{Figure scope:} The 9--11 example explains fitting; the network calculation estimates three angles from eight readings.\par
\textbf{Editable figure:} diagrams/cyber\_wls.tex\par
\textbf{[CY04] GridGuard teaching materials. Standard formulation companion: four-bus DC weighted least squares, residual tests and constrained redispatch. Local reproducible teaching fixture. 21 September 2026.:} \href{../research/source_CY04.html}{\small\ttfamily research/\allowbreak{}source\_\allowbreak{}CY04.\allowbreak{}html}\par
\textbf{[CY01] Fred C. Schweppe, J. Wildes. Power System Static-State Estimation, Part I: Exact Model. IEEE Transactions on Power Apparatus and Systems, PAS-89(1), pp. 120--125. January 1970.:} \href{https://doi.org/10.1109/TPAS.1970.292678}{\small\ttfamily https:/\allowbreak{}/\allowbreak{}doi.\allowbreak{}org/\allowbreak{}10.\allowbreak{}1109/\allowbreak{}TPAS.\allowbreak{}1970.\allowbreak{}292678}\par
\par\medskip
\subsection*{43. A residual alarm checks consistency}
\textbf{Timing:} 84:00--86:00 (2 minutes)\par
\textbf{Slide content:} Nominal J = 3.92022 is below the 15.0863 alarm threshold.
The 1\% calibration assumes the correct linear model and the stated independent Gaussian noise.\par
\textbf{Teaching notes:} The statistic is J=r\textasciicircum{}T R\textasciicircum{}-1 r, where r=z-Hxhat and R is the diagonal matrix of the assumed noise variances. Under the stated linear Gaussian model, correct H and covariance, and full column rank, the minimized statistic has m-rank(H)=5 degrees of freedom. The 99th percentile is 15.0862724694; use 15.0863 on the slide. Nominal J is 3.9202216910. Alpha=0.01 is a nominal upper-tail false-alarm probability under those assumptions. It is not a 99\% chance that the data are safe, nor the probability that an attacker is present. An isolated +0.08 p.u. change on flow12 produces J=277.4974469045 and an alarm in this seeded run. That is evidence of inconsistency, not an identification of the cause or a guarantee that all coordinated errors will be detected.\par
\textbf{Figure scope:} Upper-tail consistency test with five residual degrees of freedom. A quiet alarm does not establish trustworthy data.\par
\textbf{Editable figure:} diagrams/cyber\_residual.tex\par
\textbf{[CY04] GridGuard teaching materials. Standard formulation companion: four-bus DC weighted least squares, residual tests and constrained redispatch. Local reproducible teaching fixture. 21 September 2026.:} \href{../research/source_CY04.html}{\small\ttfamily research/\allowbreak{}source\_\allowbreak{}CY04.\allowbreak{}html}\par
\textbf{[CY03] NIST/SEMATECH. e-Handbook of Statistical Methods, Section 1.3.6.7.4: Critical Values of the Chi-Square Distribution. National Institute of Standards and Technology. Accessed 18 October 2026.:} \href{https://www.itl.nist.gov/div898/handbook/eda/section3/eda3674.htm}{\small\ttfamily https:/\allowbreak{}/\allowbreak{}www.\allowbreak{}itl.\allowbreak{}nist.\allowbreak{}gov/\allowbreak{}div898/\allowbreak{}handbook/\allowbreak{}eda/\allowbreak{}section3/\allowbreak{}eda3674.\allowbreak{}htm}\par
\par\medskip
\subsection*{44. Coordinated changes can leave the residual unchanged}
\textbf{Timing:} 86:00--88:00 (2 minutes)\par
\textbf{Slide content:} If readings and the inferred state move together, their mismatch can stay the same.
The Hc result assumes a known fixed model and access to the required unprotected measurement channels.\par
\textbf{Teaching notes:} Return to the scalar illustration: shifting both readings by three moves their best common estimate by three, while the two residuals remain -1 and +1. The linear network identity is z'=z+Hc, xhat'=xhat+c and r'=r when the same full-rank H and weights are used. Here c is a shift in the inferred angle state; Hc is the corresponding change in model predictions. Liu, Ning and Reiter establish this limitation of residual-based bad-data detection. Their result has specific access and model assumptions; it is not a claim that any attacker can alter every utility state estimate. Protected measurements restrict which changes are possible. The synthetic workshop deliberately supplies a matching coordinated perturbation as a negative control, so it can test whether a residual-only workflow mistakes consistency for trust. We do not discuss targeting real measurement channels. The operator lesson is to seek independent evidence when the trust premise is unresolved, even if the residual test is quiet.\par
\textbf{Figure scope:} Defensive explanation of a classical state-estimation limitation. The scalar picture illustrates the cancellation.\par
\textbf{Editable figure:} diagrams/cyber\_stealth.tex\par
\textbf{[CY02] Yao Liu, Peng Ning, Michael K. Reiter. False Data Injection Attacks against State Estimation in Electric Power Grids. ACM Transactions on Information and System Security, 14(1), Article 13, 33 pages. May 2011.:} \href{https://doi.org/10.1145/1952982.1952995}{\small\ttfamily https:/\allowbreak{}/\allowbreak{}doi.\allowbreak{}org/\allowbreak{}10.\allowbreak{}1145/\allowbreak{}1952982.\allowbreak{}1952995}\par
\textbf{[CY04] GridGuard teaching materials. Standard formulation companion: four-bus DC weighted least squares, residual tests and constrained redispatch. Local reproducible teaching fixture. 21 September 2026.:} \href{../research/source_CY04.html}{\small\ttfamily research/\allowbreak{}source\_\allowbreak{}CY04.\allowbreak{}html}\par
\par\medskip
\subsection*{45. A quiet residual can hide a wrong inferred state}
\textbf{Timing:} 88:00--90:00 (2 minutes)\par
\textbf{Slide content:} Isolated bias raises J; the coordinated Hc change leaves J near 3.92.
The right panel uses known synthetic truth to reveal the growing state error.\par
\textbf{Teaching notes:} Read the panels together. At scale zero both cases start from the same seeded noisy readings. Increasing one isolated flow bias raises the weighted residual and eventually crosses the threshold. The coordinated Hc changes leave J essentially unchanged while their state-error norm grows. At scale one the nominal and coordinated statistics are both 3.92022 to displayed precision, but the coordinated state error is 0.0267799957 rad, versus 0.0002374314 rad for nominal noise. The left panel's symmetric-log scale accommodates small and large statistics. The right panel compares the estimate with the known state supplied by the synthetic experiment. A real operator does not have that exact ground truth as an online test input. This is a negative-control demonstration of a detector limitation, not a statistical field evaluation, attack prevalence estimate or measurement of an LLM's detection capability.\par
\textbf{Figure scope:} Actual L8 sweep on the four-bus fixture. The left axis uses a symmetric-log scale; the right panel is an offline audit.\par
\textbf{Figure asset:} images/classroom\_cyber.png\par
\textbf{[CY05] GridGuard teaching materials. Classroom Labs L8--L9: cyber consistency sweep and trusted response gate. Local workshop notebook and companion lab functions. Reviewed 18 October 2026.:} \href{../research/source_CY05.html}{\small\ttfamily research/\allowbreak{}source\_\allowbreak{}CY05.\allowbreak{}html}\par
\textbf{[CY04] GridGuard teaching materials. Standard formulation companion: four-bus DC weighted least squares, residual tests and constrained redispatch. Local reproducible teaching fixture. 21 September 2026.:} \href{../research/source_CY04.html}{\small\ttfamily research/\allowbreak{}source\_\allowbreak{}CY04.\allowbreak{}html}\par
\textbf{[CY02] Yao Liu, Peng Ning, Michael K. Reiter. False Data Injection Attacks against State Estimation in Electric Power Grids. ACM Transactions on Information and System Security, 14(1), Article 13, 33 pages. May 2011.:} \href{https://doi.org/10.1145/1952982.1952995}{\small\ttfamily https:/\allowbreak{}/\allowbreak{}doi.\allowbreak{}org/\allowbreak{}10.\allowbreak{}1145/\allowbreak{}1952982.\allowbreak{}1952995}\par
\par\medskip
\subsection*{46. Independent trusted readings change the evidence}
\textbf{Timing:} 90:00--92:00 (2 minutes)\par
\textbf{Slide content:} The fixture supplies three separately trusted angle readings; authentication is assumed.
Three readings for three angles leave no residual redundancy. J = 0 is not a new detection result.\par
\textbf{Teaching notes:} The protected recovery uses H\_T=I\_3 with three angle readings. Each has assumed noise standard deviation 0.00005 rad and a separate seeded noise realization. Because the measurement count and rank are both three, the fit is exactly determined: residual degrees of freedom are zero, the residual is zero up to numerical precision, and the code supplies no chi-square threshold or alarm. The trust argument comes from an external assumption about these channels, not their zero residual. The workshop does not implement sensor authentication, prove channel independence operationally, or ask an LLM to decide authenticity. The trusted\_response wrapper returns HOLD when the protected set is unavailable. With the set available, it can return a numerical proposal, still with dispatch=false. In an operational design, the utility would need to justify provenance, independence, time alignment and applicability before using a replacement evidence set.\par
\textbf{Figure scope:} Exact trusted-recovery assumption: three independent angle channels, with no redundant residual test.\par
\textbf{Editable figure:} diagrams/cyber\_trusted.tex\par
\textbf{[CY04] GridGuard teaching materials. Standard formulation companion: four-bus DC weighted least squares, residual tests and constrained redispatch. Local reproducible teaching fixture. 21 September 2026.:} \href{../research/source_CY04.html}{\small\ttfamily research/\allowbreak{}source\_\allowbreak{}CY04.\allowbreak{}html}\par
\textbf{[CY05] GridGuard teaching materials. Classroom Labs L8--L9: cyber consistency sweep and trusted response gate. Local workshop notebook and companion lab functions. Reviewed 18 October 2026.:} \href{../research/source_CY05.html}{\small\ttfamily research/\allowbreak{}source\_\allowbreak{}CY05.\allowbreak{}html}\par
\par\medskip
\subsection*{47. A small QP proposes a limited correction}
\textbf{Timing:} 92:00--94:30 (2.5 minutes)\par
\textbf{Slide content:} The quadratic program (QP) finds the smallest injection satisfying every predicted branch limit with a 0.01 p.u. margin.
The trusted estimate gives u = 0.121694 p.u.; predicted flow12 becomes 0.170000 p.u.\par
\textbf{Teaching notes:} The optimizer chooses one nonnegative flexible injection u at bus 2, balanced by bus 0. The objective 0.5u\textasciicircum{}2 favors a smaller action. It is a quadratic program because this objective is quadratic while the flow constraints and action bounds are linear. Bounds are 0<=u<=0.5 p.u. Every predicted branch flow must lie between minus and plus its rating less a 0.01 p.u. planning margin. Ratings are [0.35,0.18,0.35,0.35] p.u. for branches 01,12,23,03. Flow changes use g=F B\textasciicircum{}-1 e2, with g12=-4/7, so increasing u reduces the positive 1-to-2 flow. For this protected noisy estimate the binding constraint is predicted flow12=0.17, producing u=0.1216936673 and objective 0.0074046743. This differs from the 0.1225 noiseless arithmetic example in the longer chapter. The code solves the small QP with SLSQP and independently rechecks finite values, action bounds and all predicted constraints. A failed solve or check produces HOLD.\par
\textbf{Figure scope:} Scalar convex QP in the fixed DC model. The injection at bus 2 is balanced by reference bus 0.\par
\textbf{Editable figure:} diagrams/cyber\_qp.tex\par
\textbf{[CY04] GridGuard teaching materials. Standard formulation companion: four-bus DC weighted least squares, residual tests and constrained redispatch. Local reproducible teaching fixture. 21 September 2026.:} \href{../research/source_CY04.html}{\small\ttfamily research/\allowbreak{}source\_\allowbreak{}CY04.\allowbreak{}html}\par
\par\medskip
\subsection*{48. Check the consequences on every branch}
\textbf{Timing:} 94:30--96:30 (2 minutes)\par
\textbf{Slide content:} Using the compromised estimate selects no correction; actual toy flow12 stays at 0.240000 p.u.
Trusted recovery reduces it to 0.170461 p.u., below the 0.180000 p.u. rating.\par
\textbf{Teaching notes:} Both sets of bars are evaluated using the known synthetic state after their respective actions, not the measurement estimates themselves. The action based on the coordinated corrupt estimate is u=0, leaving actual flow12 at 0.24, above its 0.18 rating. The trusted set selects u=0.1216936673. Its predicted flow12 is 0.170000; the independent offline physical audit gives actual flow12=0.1704607616. All four actual branch flows satisfy their physical ratings in this fixed toy network. The 0.01 planning margin accommodates the small protected-noise difference in this run; it is not a general robust uncertainty guarantee. Online, the code checks predicted constraints and finite values. Exact actual-state values are retained for experimental evaluation, not assumed to be an available online measurement or an operational authorization gate. The result demonstrates the value of better evidence plus a constrained solver, without establishing a benefit from an LLM.\par
\textbf{Figure scope:} Offline physical audit of the two decisions. Legend labels identify which estimate selected each action; bars show actual toy-model flows.\par
\textbf{Figure asset:} images/workshop\_trusted\_control.png\par
\textbf{[CY06] GridGuard teaching materials. Trusted control comparison: all-branch actual-flow audit for two four-bus DC decisions. Local workshop computed figure. Reviewed 18 October 2026.:} \href{../research/source_CY06.html}{\small\ttfamily research/\allowbreak{}source\_\allowbreak{}CY06.\allowbreak{}html}\par
\textbf{[CY04] GridGuard teaching materials. Standard formulation companion: four-bus DC weighted least squares, residual tests and constrained redispatch. Local reproducible teaching fixture. 21 September 2026.:} \href{../research/source_CY04.html}{\small\ttfamily research/\allowbreak{}source\_\allowbreak{}CY04.\allowbreak{}html}\par
\par\medskip
\subsection*{49. 1. Inputs: a declared four-bus DC task}
\textbf{Timing:} 96:30--98:30 (2 minutes)\par
\textbf{Slide content:} The structured input authorizes WLS\_TRUSTED\_REDISPATCH for a four-bus DC teaching case.
Trusted sensors are assumed. The separate offline audit uses known synthetic truth.\par
\textbf{Teaching notes:} Use one coherent archived run throughout these five slides. Source: C:/Users/f00863x\allowbreak{}/OPRO/output/age\allowbreak{}ntic\_grid\_workfl\allowbreak{}ow\_tutorial/resu\allowbreak{}lts/standard\_age\allowbreak{}nt\_workflow.json\allowbreak{}, runs entry with case='cyber' and supervisor='agentic', review\_id=4a10da\allowbreak{}2aa66efaa3c1f39d\allowbreak{}ae3b67372c384dbf\allowbreak{}66a297de9ed01841\allowbreak{}f6ba9c4fd9. Exact raw model journal: results/standard\allowbreak{}\_stage\_20260921T\allowbreak{}143826Z\_42d19a61\allowbreak{}/llama\_calls.jso\allowbreak{}n, call\_id=7, context (file lines 394--401). The context contains use\_case='cyber', authorized\_tool='WLS\_TRUSTED\_REDISPATCH', scope='Four-bus DC measurement/control example with explicitly assumed trusted sensors', backend='separate\_reduced\_teaching\_model', and approved\_inputs='Supplied transparent simulation fixtures; no live utility measurements'. It is a structured specification, not a recorded natural-language request. The tool authorization is already supplied, so this run does not test discovery of the right engineering method from an ambiguous request. The second bullet explains the associated numerical fixture rather than adding fields to the model's recorded input: run\_wls\_demo supplies three protected angle channels, assumes their trust, and uses known synthetic truth to audit consequences. The historical standard-stage wrapper also uses that offline violation count in its teaching verification; it is not an operational truth oracle. Preserve this distinction when presenting the screenshot. The four-bus DC case omits AC voltage magnitudes, reactive power, protection, authentication and live actuation. An illustrative operator question may be spoken separately, but must not be presented as an archived prompt.\par
\textbf{Figure scope:} Actual Awareness response containing the structured Planning input, 21 September 2026.\par
\textbf{Figure asset:} images/cyber\_input.png\par
\textbf{[TRACE\_CYBER] GridGuard teaching materials. Recorded four-bus cyber workflow: actual Llama routing, trusted-state correction and pending human review. Local executed tutorial and preserved model-call journals. 21 September 2026.:} \href{../research/source_TRACE_CYBER.html}{\small\ttfamily research/\allowbreak{}source\_\allowbreak{}TRACE\_\allowbreak{}CYBER.\allowbreak{}html}\par
\par\medskip
\subsection*{50. 2. Agent: choose the authorized estimator}
\textbf{Timing:} 98:30--100:30 (2 minutes)\par
\textbf{Slide content:} Actual model call 7 returned WLS\_TRUSTED\_REDISPATCH.
The wrapper checks that choice against the task's authorization before running the numerical tool.\par
\textbf{Teaching notes:} Exact source: C:/Users/f00863x\allowbreak{}/OPRO/output/age\allowbreak{}ntic\_grid\_workfl\allowbreak{}ow\_tutorial/resu\allowbreak{}lts/standard\_sta\allowbreak{}ge\_20260921T1438\allowbreak{}26Z\_42d19a61/lla\allowbreak{}ma\_calls.json, call\_id=7 (Planning), raw\_response.message.content at line 432 and parsed response at lines 444--449. Exact reply: \{"choice":"WLS\_T\allowbreak{}RUSTED\_REDISPATC\allowbreak{}H","reason":"Thi\allowbreak{}s tool is authorized for the given use case and scope, and it is designed for a specific DC redispatch example with trusted measurements."\}. The model is llama3.1:8b; the record contains the raw local Ollama response, with no error. This is an actual recorded LLM call from 21 September 2026, not a newly generated reply and not an LLM invocation when the notebook is read. The request uses http://localhost:11434/api/chat, a constrained JSON response schema, temperature 0 and seed 11; implementation is src/standard\_agent\_workflow.py lines 31--67. The choices are VOLTVAR\_OPF, RESTORATION\_MILP, WLS\_TRUSTED\_REDISPATCH and REVIEW. Application code requires the selected choice to equal TOOLS[case] (lines 113--120); an incompatible choice records REVIEW and leaves subsequent numerical stages NOT\_RUN. Do not describe the JSON choice as a native SDK function-call event: the wrapper reads the choice and calls Python. Nor does this routing choice authenticate sensors, select which measurements to trust, set constraints or calculate an injection. In this experiment those premises and numerical parameters are fixed. The useful visible role is bounded tool routing with a preserved reply and an explicit scope gate.\par
\textbf{Figure scope:} Actual archived llama3.1:8b reply. Python executes the tool after checking the reply.\par
\textbf{Figure asset:} images/cyber\_route.png\par
\textbf{[TRACE\_CYBER] GridGuard teaching materials. Recorded four-bus cyber workflow: actual Llama routing, trusted-state correction and pending human review. Local executed tutorial and preserved model-call journals. 21 September 2026.:} \href{../research/source_TRACE_CYBER.html}{\small\ttfamily research/\allowbreak{}source\_\allowbreak{}TRACE\_\allowbreak{}CYBER.\allowbreak{}html}\par
\par\medskip
\subsection*{51. 3. Code: calculate the provisional correction}
\textbf{Timing:} 100:30--102:30 (2 minutes)\par
\textbf{Slide content:} The tool computes the estimate and correction; the model does not calculate the setpoint.
The wrapper records the proposal, simulation checks and verification result separately.\par
\textbf{Teaching notes:} The three displayed statements are src/standard\_agent\_workflow.py lines87--89 with leading indentation removed. The builder verifies the same Python syntax tree. run\_wls\_demo computes WLS estimates, the protected-state correction and the offline known-state audit. The next two statements extract its trusted\_recovery and control records. The original lines93--96 check rank3, PROPOSE and zero known-state toy-model violations; this historical offline check is not an online truth oracle. Protected-channel trust is assumed rather than authenticated. Both supervisors call the same method and obtain the same numerical correction. In this actual agentic run, Llama subsequently returns REVIEW and the human step remains pending. Full original function and result/journal are linked under TRACE\_CYBER.\par
\textbf{Figure scope:} Exact executable wrapper code. The stage names are records of distinct operations, not independent language models.\par
\textbf{Exact source excerpt:} src/standard\_agent\_workflow.py, lines 87--89. Source statements; line breaks added for readability.
result = run\_wls\_demo()
recovery = result['trusted\_recovery']
control = recovery['control']\par
\textbf{Source alignment:} research/code\_manifest.json records source hashes, original line ranges and syntax-tree equivalence checks. research/code/ contains the readable display fragment. The complete original file is preserved in research/evidence/src/.\par
\textbf{[TRACE\_CYBER] GridGuard teaching materials. Recorded four-bus cyber workflow: actual Llama routing, trusted-state correction and pending human review. Local executed tutorial and preserved model-call journals. 21 September 2026.:} \href{../research/source_TRACE_CYBER.html}{\small\ttfamily research/\allowbreak{}source\_\allowbreak{}TRACE\_\allowbreak{}CYBER.\allowbreak{}html}\par
\par\medskip
\subsection*{52. 4. Result: a bounded correction with evidence}
\textbf{Timing:} 102:30--104:30 (2 minutes)\par
\textbf{Slide content:} The tool proposes u = 0.121694 p.u.; predicted flow12 is 0.170000 p.u.
The offline audit gives flow12 = 0.170461 p.u. and no branch-rating violations in this fixture.\par
\textbf{Teaching notes:} Source: C:/Users/f00863x\allowbreak{}/OPRO/output/age\allowbreak{}ntic\_grid\_workfl\allowbreak{}ow\_tutorial/resu\allowbreak{}lts/standard\_age\allowbreak{}nt\_workflow.json\allowbreak{}, cyber/agentic entry, summary.control and trace Planning/Execution proposal/Simulation/Verification. Exact proposed injection is 0.12169366726990237 p.u. at bus 2, balanced at reference bus 0. The objective is 0.5u\textasciicircum{}2, giving 0.007404674326798854, with 0<=u<=0.5. Predicted flows in branch order [01,12,23,03] are [0.1302813748560\allowbreak{}5382,0.169999999\allowbreak{}99999998,-0.0461\allowbreak{}97066897314544,0\allowbreak{}.148292153979369\allowbreak{}65] p.u. Physical ratings are [0.35,0.18,0.35,0.35] p.u.; predicted constraints reserve 0.01 p.u. The independent known-state toy audit gives [0.1304607615600\allowbreak{}558,0.1704607615\allowbreak{}600558,-0.047845\allowbreak{}571170041824,0.1\allowbreak{}4784557117004185\allowbreak{}] p.u. and actual\_violation\_count=0. The wrapper's modeled checks pass. Its protected estimate is [-0.019982061329\allowbreak{}599805,-0.049924\allowbreak{}46613459283,-0.0\allowbreak{}4008931656186556\allowbreak{}] rad, rank three, residual degrees of freedom zero. The three protected angles are externally assumed trustworthy; their zero residual cannot supply a redundant detection test. The compromised-estimate baseline chooses u=0 and has known-state flow12=0.24 p.u., above 0.18. This improvement comes from the supplied trusted evidence and numerical redispatch. The matched fixed supervisor computes the exact same protected numerical result and returns MODEL\_RECOMMEND; the agentic run will instead end at REVIEW. Do not attribute a better optimum, attack-detection advantage, sensor authentication or field safety to the LLM. The displayed correction remains a provisional tool result.\par
\textbf{Figure scope:} Original injection and predicted-flow excerpt; branch order 01,12,23,03. Flow12 is the second array entry.\par
\textbf{Figure asset:} images/cyber\_result.png\par
\textbf{[TRACE\_CYBER] GridGuard teaching materials. Recorded four-bus cyber workflow: actual Llama routing, trusted-state correction and pending human review. Local executed tutorial and preserved model-call journals. 21 September 2026.:} \href{../research/source_TRACE_CYBER.html}{\small\ttfamily research/\allowbreak{}source\_\allowbreak{}TRACE\_\allowbreak{}CYBER.\allowbreak{}html}\par
\par\medskip
\subsection*{53. 5. Operator: resolve trust before proceeding}
\textbf{Timing:} 104:30--106:30 (2 minutes)\par
\textbf{Slide content:} Actual model call 8 returned REVIEW despite the passed numerical checks.
The record ends at PENDING acknowledgment. A person must decide what evidence or further study is needed.\par
\textbf{Teaching notes:} Exact source: C:/Users/f00863x\allowbreak{}/OPRO/output/age\allowbreak{}ntic\_grid\_workfl\allowbreak{}ow\_tutorial/resu\allowbreak{}lts/standard\_sta\allowbreak{}ge\_20260921T1438\allowbreak{}26Z\_42d19a61/lla\allowbreak{}ma\_calls.json, call\_id=8 (Explanation), context numerical\_summary and passed\_modeled\_checks=true, raw\_response.message.content, and parsed response through line 538. The actual reply is choice='REVIEW'; its reason says the numerical workflow completed successfully and passed its modeled checks, then asks to review the results. Preserve this apparently cautious or ambiguous reply rather than replacing it with an invented confident approval. In src/standard\_agent\_workflow.py lines 134--141, the allowed explanation choices are REPORT\_LIMITS and REVIEW; REVIEW sets accepted=False. Therefore this run's final status is REVIEW even though its numerical Verification stage is PASS. This is a workflow/response-label outcome, not proof that the numerical plan is infeasible. Operator acknowledgment (lines 150--151) is PENDING with acknowledged=false, dispatch\_authorized=false and equipment\_operated=false. The review\_id is 4a10da2aa66efaa3\allowbreak{}c1f39dae3b67372c\allowbreak{}384dbf66a297de9e\allowbreak{}d01841f6ba9c4fd9\allowbreak{}. No recorded human decision follows it. The second bullet is a proposed next human step, not an event asserted to have happened. Suggested review questions: who established the protected-channel premise, whether topology and limits are current, and what AC/protection/permission checks a real application would still require. This archived run makes two actual model calls, numbered 7 and 8 within the eight-call suite; simply viewing its executed notebook makes zero new LLM calls. The numerical method and parameters match the fixed comparator, so the record establishes no superior agentic control result or reduced operator workload.\par
\textbf{Figure scope:} Actual archived explanation. No human approval, dispatch authorization or equipment operation is recorded.\par
\textbf{Figure asset:} images/cyber\_review.png\par
\textbf{[TRACE\_CYBER] GridGuard teaching materials. Recorded four-bus cyber workflow: actual Llama routing, trusted-state correction and pending human review. Local executed tutorial and preserved model-call journals. 21 September 2026.:} \href{../research/source_TRACE_CYBER.html}{\small\ttfamily research/\allowbreak{}source\_\allowbreak{}TRACE\_\allowbreak{}CYBER.\allowbreak{}html}\par
\par\medskip
\subsection*{54. Exercise: HOLD or propose for operator review?}
\textbf{Timing:} 106:30--110:30 (4 minutes)\par
\textbf{Slide content:} A: J = 3.92022; measurement trust unresolved; no independent trusted replacement.
B: J = 277.497; no trustworthy replacement is available.
C: independent trust assumed; current model; all QP checks pass; u = 0.121694. Choose the next step and name one reason.\par
\textbf{Teaching notes:} Give participants one minute to read, two minutes to compare decisions with a partner, and one minute to prepare a reason and a next evidence request. Correct answers: A is HOLD in the stated stealth-capable trust scenario; low J alone does not resolve provenance. B is HOLD and investigate the inconsistency; the statistic does not identify an attacker or authorize removal of a channel. C may be PROPOSE for operator review within this fixed DC fixture, because the trust premise is supplied externally and the current-model QP and post-solve checks pass. It is not EXECUTE. For C, all four predicted branch constraints, action bounds and finite checks matter. Expected follow-up questions include whether data are current and aligned, who established trust, whether topology or limits changed, and what operational checks the DC fixture omits. Do not reveal the answers while pairs are discussing.\par
\textbf{Figure scope:} Work in pairs. Decide whether the evidence supports progressing an advisory study, not field dispatch.\par
\textbf{[CY04] GridGuard teaching materials. Standard formulation companion: four-bus DC weighted least squares, residual tests and constrained redispatch. Local reproducible teaching fixture. 21 September 2026.:} \href{../research/source_CY04.html}{\small\ttfamily research/\allowbreak{}source\_\allowbreak{}CY04.\allowbreak{}html}\par
\textbf{[CY05] GridGuard teaching materials. Classroom Labs L8--L9: cyber consistency sweep and trusted response gate. Local workshop notebook and companion lab functions. Reviewed 18 October 2026.:} \href{../research/source_CY05.html}{\small\ttfamily research/\allowbreak{}source\_\allowbreak{}CY05.\allowbreak{}html}\par
\par\medskip
\subsection*{55. Debrief: consistency, trust and permission differ}
\textbf{Timing:} 110:30--112:30 (2 minutes)\par
\textbf{Slide content:} A and B: HOLD. A quiet residual and an alarm give different evidence, but neither repairs missing trust.
C: propose for review. Passing numerical checks does not authorize an action.
The recovery comes from trusted readings and the constrained solver; no LLM authenticates a sensor.\par
\textbf{Teaching notes:} Ask one pair to justify each packet. A passes a residual consistency test but has an unresolved trust premise, so the trusted\_response policy holds. B fails the residual test, also without a trustworthy replacement, and holds while the discrepancy is investigated. C can progress to a human review packet after the stated current-model checks. The independent trust premise still needs a real operational justification outside this demonstration. Reinforce three separate questions: do the data fit the model; why should the channels be believed; and who is permitted to act? The three trusted angles have no residual redundancy, so their zero residual supplies no extra authenticity evidence. The comparison does not isolate an LLM detection benefit or prove operational cyber resilience. A useful agent can organize the assessment and report what a numerical tool actually checked, while failing closed when necessary evidence is absent.\par
\textbf{[CY04] GridGuard teaching materials. Standard formulation companion: four-bus DC weighted least squares, residual tests and constrained redispatch. Local reproducible teaching fixture. 21 September 2026.:} \href{../research/source_CY04.html}{\small\ttfamily research/\allowbreak{}source\_\allowbreak{}CY04.\allowbreak{}html}\par
\textbf{[CY05] GridGuard teaching materials. Classroom Labs L8--L9: cyber consistency sweep and trusted response gate. Local workshop notebook and companion lab functions. Reviewed 18 October 2026.:} \href{../research/source_CY05.html}{\small\ttfamily research/\allowbreak{}source\_\allowbreak{}CY05.\allowbreak{}html}\par
\textbf{[CY02] Yao Liu, Peng Ning, Michael K. Reiter. False Data Injection Attacks against State Estimation in Electric Power Grids. ACM Transactions on Information and System Security, 14(1), Article 13, 33 pages. May 2011.:} \href{https://doi.org/10.1145/1952982.1952995}{\small\ttfamily https:/\allowbreak{}/\allowbreak{}doi.\allowbreak{}org/\allowbreak{}10.\allowbreak{}1145/\allowbreak{}1952982.\allowbreak{}1952995}\par
\par\medskip
\subsection*{56. Give the operator an evidence packet}
\textbf{Timing:} 112:30--114:00 (1.5 minutes)\par
\textbf{Slide content:} Show the trust premise, current model, proposed action, checked limits and remaining uncertainty together.
The fixture ends at PROPOSE with dispatch = false. Operator review is a separate decision.\par
\textbf{Teaching notes:} The final packet should let the operator trace a recommendation back to its evidence, rather than accept a confident explanation. Include the current reading and topology versions, who established the trusted-channel premise, the fixed DC scope, u=0.121694, predicted flow12=0.170000, all other predicted branch checks, physical ratings, 0.01 planning margin and any failed or omitted checks. State that the offline 0.170461 actual-flow audit used synthetic truth and is not an online truth oracle. The local trusted\_response function sets dispatch=false whether it holds or proposes. Any real deployment would require the utility's authorization, operating procedures and applicable additional engineering checks; these are outside this module. The agent's useful role here is to assemble the correct tool calls, preserve the evidence context, flag missing premises and explain verified results. It must not convert low residuals into trust or numerical feasibility into permission.\par
\textbf{Figure scope:} Proposed operator handoff. The classroom response is advisory and has no live control connection.\par
\textbf{Editable figure:} diagrams/cyber\_handoff.tex\par
\textbf{[CY05] GridGuard teaching materials. Classroom Labs L8--L9: cyber consistency sweep and trusted response gate. Local workshop notebook and companion lab functions. Reviewed 18 October 2026.:} \href{../research/source_CY05.html}{\small\ttfamily research/\allowbreak{}source\_\allowbreak{}CY05.\allowbreak{}html}\par
\textbf{[CY07] GridGuard teaching materials. Workshop Instructor Guide and Cyber Teaching Chapter: evidence scope and operator interpretation. Local instructor material. Reviewed 18 October 2026.:} \href{../research/source_CY07.html}{\small\ttfamily research/\allowbreak{}source\_\allowbreak{}CY07.\allowbreak{}html}\par
\par\medskip
\subsection*{57. Each failure points to a different repair}
\textbf{Timing:} 114:00--116:00 (2 minutes)\par
\textbf{Slide content:} Voltage needs an adequate model. Restoration needs the current network and service policy.
Cyber response needs justified independent evidence. A new explanation alone repairs none of these.\par
\textbf{Teaching notes:} Return to the opening comparison. A broader AC violation requires a candidate change, broader model study, or both. A changed network or rating invalidates restoration evidence until the affected study is repeated. A stealth measurement change requires an independent trust argument or escalation; a quiet residual does not supply it. Each repair is a concrete engineering step. The LLM can help interpret a request, assemble the right tool arguments and explain verified results, but it cannot substitute wording for missing evidence. Ask for one example from each case without adding another timed exercise.\par
\textbf{Figure scope:} Match the remedy to the failed evidence, rather than asking for a more confident explanation.\par
\textbf{Editable figure:} diagrams/common\_repairs.tex\par
\textbf{[W00] GridGuard project teaching materials. Classroom Labs: voltage support, six-node restoration and four-bus cyber response. Existing local GridGuard workshop; recorded numerical examples and source code. 2026; inspected 18 October 2026.:} \href{../research/workshop_evidence.md}{\small\ttfamily research/\allowbreak{}workshop\_\allowbreak{}evidence.\allowbreak{}md}\par
\textbf{[OT01] Keith Stouffer, Michael Pease, CheeYee Tang, Timothy Zimmerman, Victoria Pillitteri, Suzanne Lightman, Adam Hahn, Stephanie Saravia, Aslam Sherule, Michael Thompson. Guide to Operational Technology (OT) Security. National Institute of Standards and Technology, NIST SP 800-82 Revision 3. September 2023 (final).:} \href{https://doi.org/10.6028/NIST.SP.800-82r3}{\small\ttfamily https:/\allowbreak{}/\allowbreak{}doi.\allowbreak{}org/\allowbreak{}10.\allowbreak{}6028/\allowbreak{}NIST.\allowbreak{}SP.\allowbreak{}800-\allowbreak{}82r3}\par
\par\medskip
\subsection*{58. An operator reviews one complete recommendation packet}
\textbf{Timing:} 116:00--118:00 (2 minutes)\par
\textbf{Slide content:} Show the proposed action, expected benefit, checked limits, failed checks and remaining uncertainty.
Bind the reviewer's readback to the same case, candidate and evidence version.\par
\textbf{Teaching notes:} A useful packet lets an operator inspect why a candidate is eligible for review. It includes what the solver actually checked and what remains outside the model. Voltage should show model scope and margins. Restoration should identify the served and unserved demand and the need for operational isolation/protection work. Cyber should identify the trusted evidence assumption and the response audit. The classroom acknowledgment checks a packet identity and operator name. That local interaction has no connection to live SCADA or DMS control. In deployment, the utility's operating procedures, role permissions, interlocks and field coordination remain separate engineering requirements.\par
\textbf{Figure scope:} Illustrative advisory handoff. Local acknowledgment records review and does not dispatch equipment.\par
\textbf{Editable figure:} diagrams/common\_review\_packet.tex\par
\textbf{[W00] GridGuard project teaching materials. Classroom Labs: voltage support, six-node restoration and four-bus cyber response. Existing local GridGuard workshop; recorded numerical examples and source code. 2026; inspected 18 October 2026.:} \href{../research/workshop_evidence.md}{\small\ttfamily research/\allowbreak{}workshop\_\allowbreak{}evidence.\allowbreak{}md}\par
\textbf{[OT01] Keith Stouffer, Michael Pease, CheeYee Tang, Timothy Zimmerman, Victoria Pillitteri, Suzanne Lightman, Adam Hahn, Stephanie Saravia, Aslam Sherule, Michael Thompson. Guide to Operational Technology (OT) Security. National Institute of Standards and Technology, NIST SP 800-82 Revision 3. September 2023 (final).:} \href{https://doi.org/10.6028/NIST.SP.800-82r3}{\small\ttfamily https:/\allowbreak{}/\allowbreak{}doi.\allowbreak{}org/\allowbreak{}10.\allowbreak{}6028/\allowbreak{}NIST.\allowbreak{}SP.\allowbreak{}800-\allowbreak{}82r3}\par
\par\medskip
\subsection*{59. A useful pilot measures completed work and its errors}
\textbf{Timing:} 118:00--120:00 (2 minutes)\par
\textbf{Slide content:} Compare against the same numerical tools and a strong fixed workflow, under matched task and time budgets.
Measure correct completed tasks, missed limits, stale evidence, reviewer effort and total elapsed time.
Retain failures, input versions and decisions so another engineer can reproduce the finding.\par
\textbf{Teaching notes:} Close with a testable motivation for agentic assistance: reducing specific workflow errors or handling varied requests without weakening numerical checks. The recorded voltage selector tied voting, so the tutorial makes no blanket claim that a learned selector improves every task. The restoration improvement in critical service came from objective weights, and the cyber recovery benefit came from protected measurements. A pilot should isolate what the agent adds under matched conditions. If the question concerns human workload, use a human study instead of inferring workload from solver seconds. The teaching guide contains all three exercise debriefs, full primary references, numerical provenance and local evidence paths.\par
\textbf{[W00] GridGuard project teaching materials. Classroom Labs: voltage support, six-node restoration and four-bus cyber response. Existing local GridGuard workshop; recorded numerical examples and source code. 2026; inspected 18 October 2026.:} \href{../research/workshop_evidence.md}{\small\ttfamily research/\allowbreak{}workshop\_\allowbreak{}evidence.\allowbreak{}md}\par
\textbf{[OT01] Keith Stouffer, Michael Pease, CheeYee Tang, Timothy Zimmerman, Victoria Pillitteri, Suzanne Lightman, Adam Hahn, Stephanie Saravia, Aslam Sherule, Michael Thompson. Guide to Operational Technology (OT) Security. National Institute of Standards and Technology, NIST SP 800-82 Revision 3. September 2023 (final).:} \href{https://doi.org/10.6028/NIST.SP.800-82r3}{\small\ttfamily https:/\allowbreak{}/\allowbreak{}doi.\allowbreak{}org/\allowbreak{}10.\allowbreak{}6028/\allowbreak{}NIST.\allowbreak{}SP.\allowbreak{}800-\allowbreak{}82r3}\par
\textbf{[EVAL01] National Institute of Standards and Technology. Artificial Intelligence Risk Management Framework (AI RMF 1.0). NIST AI 100-1. 26 January 2023.:} \href{https://doi.org/10.6028/NIST.AI.100-1}{\small\ttfamily https:/\allowbreak{}/\allowbreak{}doi.\allowbreak{}org/\allowbreak{}10.\allowbreak{}6028/\allowbreak{}NIST.\allowbreak{}AI.\allowbreak{}100-\allowbreak{}1}\par
\par\medskip
\subsection*{[B01] IEEE Communications Society: official resources and network presentation template}
\textbf{Authors / issuer:} IEEE Communications Society\par
\textbf{Publisher / venue:} IEEE Communications Society\par
\textbf{Date:} Reviewed 17 October 2026\par
\textbf{Theme assets:} comsoc-lockup.png and image7.png; independent educational adaptation\par
\textbf{Full source URL:} \href{https://www.comsoc.org/about/brand-experience/resources}{\small\ttfamily https:/\allowbreak{}/\allowbreak{}www.\allowbreak{}comsoc.\allowbreak{}org/\allowbreak{}about/\allowbreak{}brand-\allowbreak{}experience/\allowbreak{}resources}\par
\par\medskip
\subsection*{[W00] Classroom Labs: voltage support, six-node restoration and four-bus cyber response}
\textbf{Authors / issuer:} GridGuard project teaching materials\par
\textbf{Publisher / venue:} Existing local GridGuard workshop; recorded numerical examples and source code\par
\textbf{Date:} 2026; inspected 18 October 2026\par
\textbf{Type:} Local reproducible teaching evidence, not a peer-reviewed deployment study\par
\textbf{Full source URL:} \href{../research/workshop_evidence.md}{\small\ttfamily research/\allowbreak{}workshop\_\allowbreak{}evidence.\allowbreak{}md}\par
\par\medskip
\subsection*{[DEF03] Building effective agents}
\textbf{Authors / issuer:} Erik S., Barry Zhang\par
\textbf{Publisher / venue:} Anthropic Engineering\par
\textbf{Date:} 2024-12-19\par
\textbf{Type:} Official practitioner architecture account by the system maintainers\par
\textbf{Location:} What are agents?; Agents; Workflow: Routing; Workflow: Orchestrator-workers; Workflow: Evaluator-optimizer\par
\textbf{Version note:} The live page now warns that tooling has changed since the original publication. Use its conceptual distinction, not its present SDK list as a claim about the 2024 implementation.\par
\textbf{Full source URL:} \href{https://www.anthropic.com/engineering/building-effective-agents}{\small\ttfamily https:/\allowbreak{}/\allowbreak{}www.\allowbreak{}anthropic.\allowbreak{}com/\allowbreak{}engineering/\allowbreak{}building-\allowbreak{}effective-\allowbreak{}agents}\par
\par\medskip
\subsection*{[OT01] Guide to Operational Technology (OT) Security}
\textbf{Authors / issuer:} Keith Stouffer, Michael Pease, CheeYee Tang, Timothy Zimmerman, Victoria Pillitteri, Suzanne Lightman, Adam Hahn, Stephanie Saravia, Aslam Sherule, Michael Thompson\par
\textbf{Publisher / venue:} National Institute of Standards and Technology, NIST SP 800-82 Revision 3\par
\textbf{Date:} September 2023 (final)\par
\textbf{Type:} Published OT security guidance\par
\textbf{Full source URL:} \href{https://doi.org/10.6028/NIST.SP.800-82r3}{\small\ttfamily https:/\allowbreak{}/\allowbreak{}doi.\allowbreak{}org/\allowbreak{}10.\allowbreak{}6028/\allowbreak{}NIST.\allowbreak{}SP.\allowbreak{}800-\allowbreak{}82r3}\par
\textbf{Full paper PDF:} \href{https://nvlpubs.nist.gov/nistpubs/SpecialPublications/NIST.SP.800-82r3.pdf}{\small\ttfamily https:/\allowbreak{}/\allowbreak{}nvlpubs.\allowbreak{}nist.\allowbreak{}gov/\allowbreak{}nistpubs/\allowbreak{}SpecialPublications/\allowbreak{}NIST.\allowbreak{}SP.\allowbreak{}800-\allowbreak{}82r3.\allowbreak{}pdf}\par
\textbf{Publication:} \href{https://csrc.nist.gov/pubs/sp/800/82/r3/final}{\small\ttfamily https:/\allowbreak{}/\allowbreak{}csrc.\allowbreak{}nist.\allowbreak{}gov/\allowbreak{}pubs/\allowbreak{}sp/\allowbreak{}800/\allowbreak{}82/\allowbreak{}r3/\allowbreak{}final}\par
\par\medskip
\subsection*{[VCASE] Synthetic Greensboro feeder uhs20\_1247--udt8894: model map and declared study scope}
\textbf{Authors / issuer:} Existing GridGuard workshop / SMART-DS synthetic-feeder extraction\par
\textbf{Publisher / venue:} data.openei.org\par
\textbf{Date:} Existing teaching snapshot reviewed 18 October 2026\par
\textbf{Type:} Existing local scientific model figure; external URL is the dataset reference retained by the original workshop\par
\textbf{Full source URL:} \href{https://data.openei.org/submissions/2981}{\small\ttfamily https:/\allowbreak{}/\allowbreak{}data.\allowbreak{}openei.\allowbreak{}org/\allowbreak{}submissions/\allowbreak{}2981}\par
\par\medskip
\subsection*{[VGUIDE] Instructor guide: a coherent 180-minute teaching session}
\textbf{Authors / issuer:} GridGuard teaching materials\par
\textbf{Publisher / venue:} GridGuard project / local snapshot\par
\textbf{Date:} Existing local workshop snapshot, reviewed 18 October 2026\par
\textbf{Type:} Local instructor guide\par
\textbf{Full source URL:} \href{../research/source_VGUIDE.html}{\small\ttfamily research/\allowbreak{}source\_\allowbreak{}VGUIDE.\allowbreak{}html}\par
\par\medskip
\subsection*{[VCTRL] Continuous fleet P/Q control and classroom validation/readback helpers}
\textbf{Authors / issuer:} GridGuard teaching materials\par
\textbf{Publisher / venue:} GridGuard project / local snapshot\par
\textbf{Date:} Local source reviewed 18 October 2026\par
\textbf{Type:} Actual local numerical/control and review implementation\par
\textbf{Full source URL:} \href{../research/source_VCTRL.html}{\small\ttfamily research/\allowbreak{}source\_\allowbreak{}VCTRL.\allowbreak{}html}\par
\par\medskip
\subsection*{[VNOTE] Agentic AI for grid operators --- Classroom Labs}
\textbf{Authors / issuer:} GridGuard teaching materials\par
\textbf{Publisher / venue:} GridGuard project / local snapshot\par
\textbf{Date:} Executed 17 October 2026\par
\textbf{Type:} Executed notebook with stored outputs\par
\textbf{Full source URL:} \href{../research/source_VNOTE.html}{\small\ttfamily research/\allowbreak{}source\_\allowbreak{}VNOTE.\allowbreak{}html}\par
\par\medskip
\subsection*{[VMPC] Model Predictive Control: Theory, Computation, and Design, second edition}
\textbf{Authors / issuer:} James B. Rawlings; David Q. Mayne; Moritz M. Diehl\par
\textbf{Publisher / venue:} Nob Hill Publishing\par
\textbf{Date:} 2017\par
\textbf{Type:} Primary author-hosted textbook\par
\textbf{Full source URL:} \href{https://sites.engineering.ucsb.edu/~jbraw/mpc/MPC-book-2nd-edition-1st-printing.pdf}{\small\ttfamily https:/\allowbreak{}/\allowbreak{}sites.\allowbreak{}engineering.\allowbreak{}ucsb.\allowbreak{}edu/\allowbreak{}\textasciitilde{}jbraw/\allowbreak{}mpc/\allowbreak{}MPC-\allowbreak{}book-\allowbreak{}2nd-\allowbreak{}edition-\allowbreak{}1st-\allowbreak{}printing.\allowbreak{}pdf}\par
\par\medskip
\subsection*{[VSAC] Soft Actor-Critic: Off-Policy Maximum Entropy Deep Reinforcement Learning with a Stochastic Actor}
\textbf{Authors / issuer:} Tuomas Haarnoja; Aurick Zhou; Pieter Abbeel; Sergey Levine\par
\textbf{Publisher / venue:} Proceedings of the 35th International Conference on Machine Learning, PMLR 80, pp. 1861--1870\par
\textbf{Date:} 2018\par
\textbf{Type:} Primary research paper\par
\textbf{Full source URL:} \href{https://proceedings.mlr.press/v80/haarnoja18b.html}{\small\ttfamily https:/\allowbreak{}/\allowbreak{}proceedings.\allowbreak{}mlr.\allowbreak{}press/\allowbreak{}v80/\allowbreak{}haarnoja18b.\allowbreak{}html}\par
\par\medskip
\subsection*{[VPILOT] Voltage response-selection pilot: how to read the evidence}
\textbf{Authors / issuer:} GridGuard teaching materials\par
\textbf{Publisher / venue:} GridGuard project / local snapshot\par
\textbf{Date:} Protocol frozen 27 September 2026; reviewed 18 October 2026\par
\textbf{Type:} Local experiment protocol and evidence guide\par
\textbf{Full source URL:} \href{../research/source_VPILOT.html}{\small\ttfamily research/\allowbreak{}source\_\allowbreak{}VPILOT.\allowbreak{}html}\par
\par\medskip
\subsection*{[VARCH] Archived immutable voltage contexts and actual local Llama/Qwen response records}
\textbf{Authors / issuer:} GridGuard teaching materials\par
\textbf{Publisher / venue:} GridGuard project / local snapshot\par
\textbf{Date:} Recorded September 2026; reviewed 18 October 2026\par
\textbf{Type:} Original local numerical contexts and actual model-response journals\par
\textbf{Full source URL:} \href{../research/source_VARCH.html}{\small\ttfamily research/\allowbreak{}source\_\allowbreak{}VARCH.\allowbreak{}html}\par
\par\medskip
\subsection*{[VRECHECK] Distinct final first-action numerical rechecks after bounded response selection}
\textbf{Authors / issuer:} GridGuard teaching materials\par
\textbf{Publisher / venue:} GridGuard project / local snapshot\par
\textbf{Date:} Recorded September 2026\par
\textbf{Type:} Original saved numerical results\par
\textbf{Full source URL:} \href{../research/source_VRECHECK.html}{\small\ttfamily research/\allowbreak{}source\_\allowbreak{}VRECHECK.\allowbreak{}html}\par
\par\medskip
\subsection*{[TRACE\_VOLT] Recorded OpenDSS voltage workflow: P3 case, actual Llama replies, dashboard and Python source}
\textbf{Authors / issuer:} GridGuard teaching materials\par
\textbf{Publisher / venue:} Local executed tutorial and preserved model-call journals\par
\textbf{Date:} 27 September 2026\par
\textbf{Type:} Archived local demonstration, not a field deployment\par
\textbf{Full source URL:} \href{../research/source_TRACE_VOLT.html}{\small\ttfamily research/\allowbreak{}source\_\allowbreak{}TRACE\_\allowbreak{}VOLT.\allowbreak{}html}\par
\par\medskip
\subsection*{[VANAL] Frozen response-selector comparison: held-out MV test and paired full-feeder AC transfer}
\textbf{Authors / issuer:} GridGuard teaching materials\par
\textbf{Publisher / venue:} GridGuard project / local snapshot\par
\textbf{Date:} Protocol frozen 27 September 2026\par
\textbf{Type:} Original local analysis of recorded outcomes\par
\textbf{Full source URL:} \href{../research/source_VANAL.html}{\small\ttfamily research/\allowbreak{}source\_\allowbreak{}VANAL.\allowbreak{}html}\par
\par\medskip
\subsection*{[VAC] Paired same-action full-feeder audit and independent OpenDSS MPC first-action plans}
\textbf{Authors / issuer:} GridGuard teaching materials\par
\textbf{Publisher / venue:} GridGuard project / local snapshot\par
\textbf{Date:} Recorded September 2026\par
\textbf{Type:} Original saved full-model comparison\par
\textbf{Full source URL:} \href{../research/source_VAC.html}{\small\ttfamily research/\allowbreak{}source\_\allowbreak{}VAC.\allowbreak{}html}\par
\par\medskip
\subsection*{[RST01] Agentic AI for Grid Operators: restoration teaching chapter, frames R01--R18}
\textbf{Authors / issuer:} GridGuard teaching materials\par
\textbf{Publisher / venue:} Local existing workshop teaching edition\par
\textbf{Date:} 17 October 2026 edition\par
\textbf{Full source URL:} \href{../research/source_RST01.html}{\small\ttfamily research/\allowbreak{}source\_\allowbreak{}RST01.\allowbreak{}html}\par
\par\medskip
\subsection*{[RST02] Classroom Labs: L6 restoration sweep and L7 changed-contract binding}
\textbf{Authors / issuer:} GridGuard teaching materials\par
\textbf{Publisher / venue:} Local executed Jupyter notebook\par
\textbf{Date:} 17 October 2026 teaching edition; inspected18 October2026\par
\textbf{Full source URL:} \href{../research/source_RST02.html}{\small\ttfamily research/\allowbreak{}source\_\allowbreak{}RST02.\allowbreak{}html}\par
\par\medskip
\subsection*{[RST03] Six-node restoration graph, MILP and exhaustive enumeration implementation}
\textbf{Authors / issuer:} GridGuard teaching materials\par
\textbf{Publisher / venue:} Local Python source and saved numerical evidence\par
\textbf{Date:} Existing17 October2026 teaching edition; inspected18 October2026\par
\textbf{Full source URL:} \href{../research/source_RST03.html}{\small\ttfamily research/\allowbreak{}source\_\allowbreak{}RST03.\allowbreak{}html}\par
\par\medskip
\subsection*{[RSTL01] On the Radiality Constraints for Distribution System Restoration and Reconfiguration Problems}
\textbf{Authors / issuer:} Ying Wang, Yin Xu, Jiaxu Li, Jinghan He, Xiaojun Wang\par
\textbf{Publisher / venue:} Author preprint, arXiv:1912.05185v3\par
\textbf{Date:} 15 March2020(v3); first submitted11 December2019\par
\textbf{Full source URL:} \href{https://arxiv.org/abs/1912.05185v3}{\small\ttfamily https:/\allowbreak{}/\allowbreak{}arxiv.\allowbreak{}org/\allowbreak{}abs/\allowbreak{}1912.\allowbreak{}05185v3}\par
\textbf{Full paper PDF:} \href{https://arxiv.org/pdf/1912.05185v3}{\small\ttfamily https:/\allowbreak{}/\allowbreak{}arxiv.\allowbreak{}org/\allowbreak{}pdf/\allowbreak{}1912.\allowbreak{}05185v3}\par
\par\medskip
\subsection*{[RST05] Classroom restoration capacity sweep: total and critical service}
\textbf{Authors / issuer:} GridGuard teaching materials\par
\textbf{Publisher / venue:} Original computed local scientific figure\par
\textbf{Date:} 17 October2026 teaching edition; inspected18 October2026\par
\textbf{Full source URL:} \href{../research/source_RST05.html}{\small\ttfamily research/\allowbreak{}source\_\allowbreak{}RST05.\allowbreak{}html}\par
\par\medskip
\subsection*{[TRACE\_REST] Recorded six-node restoration workflow: actual Llama routing, MILP evidence and pending human review}
\textbf{Authors / issuer:} GridGuard teaching materials\par
\textbf{Publisher / venue:} Local executed tutorial and preserved model-call journals\par
\textbf{Date:} 21 September 2026\par
\textbf{Type:} Archived local demonstration, not a field deployment\par
\textbf{Full source URL:} \href{../research/source_TRACE_REST.html}{\small\ttfamily research/\allowbreak{}source\_\allowbreak{}TRACE\_\allowbreak{}REST.\allowbreak{}html}\par
\par\medskip
\subsection*{[RST04] Evidence digest, binding and contract matching functions}
\textbf{Authors / issuer:} GridGuard teaching materials\par
\textbf{Publisher / venue:} Local classroom implementation\par
\textbf{Date:} Existing17 October2026 teaching edition; inspected18 October2026\par
\textbf{Full source URL:} \href{../research/source_RST04.html}{\small\ttfamily research/\allowbreak{}source\_\allowbreak{}RST04.\allowbreak{}html}\par
\par\medskip
\subsection*{[RST06] Fresh restoration enumeration and explicit topology-binding teaching check}
\textbf{Authors / issuer:} GridGuard teaching materials\par
\textbf{Publisher / venue:} Local reproducible teaching evidence\par
\textbf{Date:} 18 October2026\par
\textbf{Full source URL:} \href{../research/source_RST06.html}{\small\ttfamily research/\allowbreak{}source\_\allowbreak{}RST06.\allowbreak{}html}\par
\par\medskip
\subsection*{[RSTL02] Sequential Service Restoration for Unbalanced Distribution Systems and Microgrids}
\textbf{Authors / issuer:} Bo Chen, Chen Chen, Jianhui Wang, Karen L. Butler-Purry\par
\textbf{Publisher / venue:} IEEE Transactions on Power Systems33(2),1507--1520\par
\textbf{Date:} March2018\par
\textbf{Full source URL:} \href{https://doi.org/10.1109/TPWRS.2017.2720122}{\small\ttfamily https:/\allowbreak{}/\allowbreak{}doi.\allowbreak{}org/\allowbreak{}10.\allowbreak{}1109/\allowbreak{}TPWRS.\allowbreak{}2017.\allowbreak{}2720122}\par
\par\medskip
\subsection*{[CY04] Standard formulation companion: four-bus DC weighted least squares, residual tests and constrained redispatch}
\textbf{Authors / issuer:} GridGuard teaching materials\par
\textbf{Publisher / venue:} Local reproducible teaching fixture\par
\textbf{Date:} Recorded 21 September 2026; reviewed 18 October 2026\par
\textbf{Full source URL:} \href{../research/source_CY04.html}{\small\ttfamily research/\allowbreak{}source\_\allowbreak{}CY04.\allowbreak{}html}\par
\par\medskip
\subsection*{[CY07] Workshop Instructor Guide and Cyber Teaching Chapter: evidence scope and operator interpretation}
\textbf{Authors / issuer:} GridGuard teaching materials\par
\textbf{Publisher / venue:} Local instructor material\par
\textbf{Date:} Reviewed 18 October 2026\par
\textbf{Full source URL:} \href{../research/source_CY07.html}{\small\ttfamily research/\allowbreak{}source\_\allowbreak{}CY07.\allowbreak{}html}\par
\par\medskip
\subsection*{[CY01] Power System Static-State Estimation, Part I: Exact Model}
\textbf{Authors / issuer:} Fred C. Schweppe, J. Wildes\par
\textbf{Publisher / venue:} IEEE Transactions on Power Apparatus and Systems, PAS-89(1), pp. 120--125\par
\textbf{Date:} January 1970\par
\textbf{Full source URL:} \href{https://doi.org/10.1109/TPAS.1970.292678}{\small\ttfamily https:/\allowbreak{}/\allowbreak{}doi.\allowbreak{}org/\allowbreak{}10.\allowbreak{}1109/\allowbreak{}TPAS.\allowbreak{}1970.\allowbreak{}292678}\par
\par\medskip
\subsection*{[CY03] e-Handbook of Statistical Methods, Section 1.3.6.7.4: Critical Values of the Chi-Square Distribution}
\textbf{Authors / issuer:} NIST/SEMATECH\par
\textbf{Publisher / venue:} National Institute of Standards and Technology\par
\textbf{Date:} Undated online reference; accessed 18 October 2026\par
\textbf{Full source URL:} \href{https://www.itl.nist.gov/div898/handbook/eda/section3/eda3674.htm}{\small\ttfamily https:/\allowbreak{}/\allowbreak{}www.\allowbreak{}itl.\allowbreak{}nist.\allowbreak{}gov/\allowbreak{}div898/\allowbreak{}handbook/\allowbreak{}eda/\allowbreak{}section3/\allowbreak{}eda3674.\allowbreak{}htm}\par
\par\medskip
\subsection*{[CY02] False Data Injection Attacks against State Estimation in Electric Power Grids}
\textbf{Authors / issuer:} Yao Liu, Peng Ning, Michael K. Reiter\par
\textbf{Publisher / venue:} ACM Transactions on Information and System Security, 14(1), Article 13, 33 pages\par
\textbf{Date:} May 2011\par
\textbf{Full source URL:} \href{https://doi.org/10.1145/1952982.1952995}{\small\ttfamily https:/\allowbreak{}/\allowbreak{}doi.\allowbreak{}org/\allowbreak{}10.\allowbreak{}1145/\allowbreak{}1952982.\allowbreak{}1952995}\par
\textbf{Full paper PDF:} \href{https://reitermk.github.io/papers/2011/TISSEC2.pdf}{\small\ttfamily https:/\allowbreak{}/\allowbreak{}reitermk.\allowbreak{}github.\allowbreak{}io/\allowbreak{}papers/\allowbreak{}2011/\allowbreak{}TISSEC2.\allowbreak{}pdf}\par
\par\medskip
\subsection*{[CY05] Classroom Labs L8--L9: cyber consistency sweep and trusted response gate}
\textbf{Authors / issuer:} GridGuard teaching materials\par
\textbf{Publisher / venue:} Local workshop notebook and companion lab functions\par
\textbf{Date:} Reviewed 18 October 2026\par
\textbf{Full source URL:} \href{../research/source_CY05.html}{\small\ttfamily research/\allowbreak{}source\_\allowbreak{}CY05.\allowbreak{}html}\par
\par\medskip
\subsection*{[CY06] Trusted control comparison: all-branch actual-flow audit for two four-bus DC decisions}
\textbf{Authors / issuer:} GridGuard teaching materials\par
\textbf{Publisher / venue:} Local workshop computed figure\par
\textbf{Date:} Reviewed 18 October 2026\par
\textbf{Full source URL:} \href{../research/source_CY06.html}{\small\ttfamily research/\allowbreak{}source\_\allowbreak{}CY06.\allowbreak{}html}\par
\par\medskip
\subsection*{[TRACE\_CYBER] Recorded four-bus cyber workflow: actual Llama routing, trusted-state correction and pending human review}
\textbf{Authors / issuer:} GridGuard teaching materials\par
\textbf{Publisher / venue:} Local executed tutorial and preserved model-call journals\par
\textbf{Date:} 21 September 2026\par
\textbf{Type:} Archived local demonstration, not a field deployment\par
\textbf{Full source URL:} \href{../research/source_TRACE_CYBER.html}{\small\ttfamily research/\allowbreak{}source\_\allowbreak{}TRACE\_\allowbreak{}CYBER.\allowbreak{}html}\par
\par\medskip
\subsection*{[EVAL01] Artificial Intelligence Risk Management Framework (AI RMF 1.0)}
\textbf{Authors / issuer:} National Institute of Standards and Technology\par
\textbf{Publisher / venue:} NIST AI 100-1\par
\textbf{Date:} 26 January 2023\par
\textbf{Type:} Voluntary risk-management guidance\par
\textbf{Full source URL:} \href{https://doi.org/10.6028/NIST.AI.100-1}{\small\ttfamily https:/\allowbreak{}/\allowbreak{}doi.\allowbreak{}org/\allowbreak{}10.\allowbreak{}6028/\allowbreak{}NIST.\allowbreak{}AI.\allowbreak{}100-\allowbreak{}1}\par
\textbf{Full paper PDF:} \href{https://nvlpubs.nist.gov/nistpubs/ai/NIST.AI.100-1.pdf}{\small\ttfamily https:/\allowbreak{}/\allowbreak{}nvlpubs.\allowbreak{}nist.\allowbreak{}gov/\allowbreak{}nistpubs/\allowbreak{}ai/\allowbreak{}NIST.\allowbreak{}AI.\allowbreak{}100-\allowbreak{}1.\allowbreak{}pdf}\par
\par\medskip
\subsection*{Visual credit: feeder.png}
\textbf{Caption:} Existing SMART-DS synthetic Greensboro feeder; one teaching model, not a live utility circuit.\par
\textbf{Credit:} Existing GridGuard workshop numerical figure\par
\textbf{Used on:} 6. Voltage support: return a command with evidence\par
\textbf{Asset type:} Synthetic feeder visualization\par
\textbf{Transformation:} Unchanged original plot.\par
\textbf{Evidence limit:} Existing SMART-DS synthetic Greensboro feeder; one teaching model, not a live utility circuit.\par
\textbf{Original source / page:} \href{../research/source_VCASE.html}{\small\ttfamily research/\allowbreak{}source\_\allowbreak{}VCASE.\allowbreak{}html}\par
\textbf{Original figure / image URL:} \href{../research/evidence/figures/feeder.png}{\small\ttfamily research/\allowbreak{}evidence/\allowbreak{}figures/\allowbreak{}feeder.\allowbreak{}png}\par
\par\medskip
\subsection*{Visual credit: voltage\_current\_comparison.png}
\textbf{Caption:} Replotted stored L2 current checks. 100\% is the declared allowed-current limit.\par
\textbf{Credit:} Original teaching replot using existing GridGuard records\par
\textbf{Used on:} 9. Compare two freshly checked proposals\par
\textbf{Asset type:} Recorded numerical teaching figure\par
\textbf{Transformation:} New larger vector chart from the stored six-decimal L2 output values; no new numerical experiment.\par
\textbf{Evidence limit:} Replotted stored L2 current checks. 100\% is the declared allowed-current limit.\par
\textbf{Original source / page:} \href{../research/source_VNOTE.html}{\small\ttfamily research/\allowbreak{}source\_\allowbreak{}VNOTE.\allowbreak{}html}\par
\textbf{Original figure / image URL:} \href{../research/evidence/replots/voltage_current_comparison.png}{\small\ttfamily research/\allowbreak{}evidence/\allowbreak{}replots/\allowbreak{}voltage\_\allowbreak{}current\_\allowbreak{}comparison.\allowbreak{}png}\par
\par\medskip
\subsection*{Visual credit: voltage\_input.png}
\textbf{Caption:} Actual dashboard Awareness excerpt, P3 OpenDSS, 27 September 2026; separate from test\_01.\par
\textbf{Credit:} GridGuard project: browser capture of archived executed notebook/dashboard\par
\textbf{Used on:} 12. 1. Inputs: the recorded OpenDSS case\par
\textbf{Asset type:} Actual local GridGuard screenshot\par
\textbf{Transformation:} Existing HTML rendered in a browser; selected element crop; text enlarged without changing recorded values. Full provenance: research/screenshots/capture\_manifest.json.\par
\textbf{Evidence limit:} Actual dashboard Awareness excerpt, P3 OpenDSS, 27 September 2026; separate from test\_01.\par
\textbf{Original source / page:} \href{../research/source_TRACE_VOLT.html}{\small\ttfamily research/\allowbreak{}source\_\allowbreak{}TRACE\_\allowbreak{}VOLT.\allowbreak{}html}\par
\textbf{Original figure / image URL:} \href{../research/evidence/additional/voltage_input.png}{\small\ttfamily research/\allowbreak{}evidence/\allowbreak{}additional/\allowbreak{}voltage\_\allowbreak{}input.\allowbreak{}png}\par
\par\medskip
\subsection*{Visual credit: voltage\_route.png}
\textbf{Caption:} Actual dashboard Recovery reply after the thermal rejection; no new model call was made for this slide.\par
\textbf{Credit:} GridGuard project: browser capture of archived executed notebook/dashboard\par
\textbf{Used on:} 13. 2. Agent: request the MPC fallback\par
\textbf{Asset type:} Actual local GridGuard screenshot\par
\textbf{Transformation:} Existing HTML rendered in a browser; selected element crop; text enlarged without changing recorded values. Full provenance: research/screenshots/capture\_manifest.json.\par
\textbf{Evidence limit:} Actual dashboard Recovery reply after the thermal rejection; no new model call was made for this slide.\par
\textbf{Original source / page:} \href{../research/source_TRACE_VOLT.html}{\small\ttfamily research/\allowbreak{}source\_\allowbreak{}TRACE\_\allowbreak{}VOLT.\allowbreak{}html}\par
\textbf{Original figure / image URL:} \href{../research/evidence/additional/voltage_route.png}{\small\ttfamily research/\allowbreak{}evidence/\allowbreak{}additional/\allowbreak{}voltage\_\allowbreak{}route.\allowbreak{}png}\par
\par\medskip
\subsection*{Visual credit: voltage\_result.png}
\textbf{Caption:} Recorded P3 first-action results. 100\% denotes the declared current allowance, which is 80\% of the nominal rating.\par
\textbf{Credit:} GridGuard project: browser capture of archived executed notebook/dashboard\par
\textbf{Used on:} 15. 4. Result: SAC fails; MPC passes\par
\textbf{Asset type:} Actual local GridGuard screenshot\par
\textbf{Transformation:} Existing HTML rendered in a browser; selected element crop; text enlarged without changing recorded values. Full provenance: research/screenshots/capture\_manifest.json.\par
\textbf{Evidence limit:} Recorded P3 first-action results. 100\% denotes the declared current allowance, which is 80\% of the nominal rating.\par
\textbf{Original source / page:} \href{../research/source_TRACE_VOLT.html}{\small\ttfamily research/\allowbreak{}source\_\allowbreak{}TRACE\_\allowbreak{}VOLT.\allowbreak{}html}\par
\textbf{Original figure / image URL:} \href{../research/evidence/additional/voltage_result.png}{\small\ttfamily research/\allowbreak{}evidence/\allowbreak{}additional/\allowbreak{}voltage\_\allowbreak{}result.\allowbreak{}png}\par
\par\medskip
\subsection*{Visual credit: voltage\_review.png}
\textbf{Caption:} Original operator controls, still awaiting human readback. No acknowledgment or equipment operation was performed.\par
\textbf{Credit:} GridGuard project: browser capture of archived executed notebook/dashboard\par
\textbf{Used on:} 16. 5. Operator: review the exact recommendation\par
\textbf{Asset type:} Actual local GridGuard screenshot\par
\textbf{Transformation:} Existing HTML rendered in a browser; selected element crop; text enlarged without changing recorded values. Full provenance: research/screenshots/capture\_manifest.json.\par
\textbf{Evidence limit:} Original operator controls, still awaiting human readback. No acknowledgment or equipment operation was performed.\par
\textbf{Original source / page:} \href{../research/source_TRACE_VOLT.html}{\small\ttfamily research/\allowbreak{}source\_\allowbreak{}TRACE\_\allowbreak{}VOLT.\allowbreak{}html}\par
\textbf{Original figure / image URL:} \href{../research/evidence/additional/voltage_review.png}{\small\ttfamily research/\allowbreak{}evidence/\allowbreak{}additional/\allowbreak{}voltage\_\allowbreak{}review.\allowbreak{}png}\par
\par\medskip
\subsection*{Visual credit: voltage\_scope\_commands.png}
\textbf{Caption:} Replotted paired command records. Each curve uses a separately planned control; this is not a solver-error plot.\par
\textbf{Credit:} Original teaching replot using existing GridGuard records\par
\textbf{Used on:} 21. Replan with the model that covers the request\par
\textbf{Asset type:} Recorded numerical teaching figure\par
\textbf{Transformation:} New larger vector chart from the paired 12-context command records; values unchanged.\par
\textbf{Evidence limit:} Replotted paired command records. Each curve uses a separately planned control; this is not a solver-error plot.\par
\textbf{Original source / page:} \href{../research/source_VAC.html}{\small\ttfamily research/\allowbreak{}source\_\allowbreak{}VAC.\allowbreak{}html}\par
\textbf{Original figure / image URL:} \href{../research/evidence/replots/voltage_scope_commands.png}{\small\ttfamily research/\allowbreak{}evidence/\allowbreak{}replots/\allowbreak{}voltage\_\allowbreak{}scope\_\allowbreak{}commands.\allowbreak{}png}\par
\par\medskip
\subsection*{Visual credit: classroom\_restoration.png}
\textbf{Caption:} Existing computed Classroom Labs L6 figure; both objectives use the same feasible network.\par
\textbf{Credit:} Existing GridGuard workshop numerical figure\par
\textbf{Used on:} 29. A higher tie rating does not always restore another load\par
\textbf{Asset type:} Recorded numerical teaching figure\par
\textbf{Transformation:} Unchanged original plot.\par
\textbf{Evidence limit:} Existing computed Classroom Labs L6 figure; both objectives use the same feasible network.\par
\textbf{Original source / page:} \href{../research/source_RST02.html}{\small\ttfamily research/\allowbreak{}source\_\allowbreak{}RST02.\allowbreak{}html}\par
\textbf{Original figure / image URL:} \href{../research/evidence/figures/classroom_restoration.png}{\small\ttfamily research/\allowbreak{}evidence/\allowbreak{}figures/\allowbreak{}classroom\_\allowbreak{}restoration.\allowbreak{}png}\par
\par\medskip
\subsection*{Visual credit: rest\_input.png}
\textbf{Caption:} Actual Awareness response containing the structured Planning input, 21 September 2026; tie0--4 = 100 kW.\par
\textbf{Credit:} GridGuard project: browser capture of archived executed notebook/dashboard\par
\textbf{Used on:} 30. 1. Inputs: the six-node restoration task\par
\textbf{Asset type:} Actual local GridGuard screenshot\par
\textbf{Transformation:} Existing HTML rendered in a browser; selected element crop; text enlarged without changing recorded values. Full provenance: research/screenshots/capture\_manifest.json.\par
\textbf{Evidence limit:} Actual Awareness response containing the structured Planning input, 21 September 2026; tie0--4 = 100 kW.\par
\textbf{Original source / page:} \href{../research/source_TRACE_REST.html}{\small\ttfamily research/\allowbreak{}source\_\allowbreak{}TRACE\_\allowbreak{}REST.\allowbreak{}html}\par
\textbf{Original figure / image URL:} \href{../research/evidence/additional/rest_input.png}{\small\ttfamily research/\allowbreak{}evidence/\allowbreak{}additional/\allowbreak{}rest\_\allowbreak{}input.\allowbreak{}png}\par
\par\medskip
\subsection*{Visual credit: rest\_route.png}
\textbf{Caption:} Crop of the original Planning response: actual Llama choice, returned reason, call ID and provenance.\par
\textbf{Credit:} GridGuard project: browser capture of archived executed notebook/dashboard\par
\textbf{Used on:} 31. 2. Agent: choose the authorized MILP\par
\textbf{Asset type:} Actual local GridGuard screenshot\par
\textbf{Transformation:} Existing HTML rendered in a browser; selected element crop; text enlarged without changing recorded values. Full provenance: research/screenshots/capture\_manifest.json.\par
\textbf{Evidence limit:} Crop of the original Planning response: actual Llama choice, returned reason, call ID and provenance.\par
\textbf{Original source / page:} \href{../research/source_TRACE_REST.html}{\small\ttfamily research/\allowbreak{}source\_\allowbreak{}TRACE\_\allowbreak{}REST.\allowbreak{}html}\par
\textbf{Original figure / image URL:} \href{../research/evidence/additional/rest_route.png}{\small\ttfamily research/\allowbreak{}evidence/\allowbreak{}additional/\allowbreak{}rest\_\allowbreak{}route.\allowbreak{}png}\par
\par\medskip
\subsection*{Visual credit: rest\_result.png}
\textbf{Caption:} Original numerical-plan excerpt. Separate verification records solver success and enumeration agreement.\par
\textbf{Credit:} GridGuard project: browser capture of archived executed notebook/dashboard\par
\textbf{Used on:} 33. 4. Result: 200 kW total, 90 kW critical\par
\textbf{Asset type:} Actual local GridGuard screenshot\par
\textbf{Transformation:} Existing HTML rendered in a browser; selected element crop; text enlarged without changing recorded values. Full provenance: research/screenshots/capture\_manifest.json.\par
\textbf{Evidence limit:} Original numerical-plan excerpt. Separate verification records solver success and enumeration agreement.\par
\textbf{Original source / page:} \href{../research/source_TRACE_REST.html}{\small\ttfamily research/\allowbreak{}source\_\allowbreak{}TRACE\_\allowbreak{}REST.\allowbreak{}html}\par
\textbf{Original figure / image URL:} \href{../research/evidence/additional/rest_result.png}{\small\ttfamily research/\allowbreak{}evidence/\allowbreak{}additional/\allowbreak{}rest\_\allowbreak{}result.\allowbreak{}png}\par
\par\medskip
\subsection*{Visual credit: rest\_review.png}
\textbf{Caption:} Original Explanation output. The separate acknowledgment record remains PENDING, with no dispatch or equipment operation.\par
\textbf{Credit:} GridGuard project: browser capture of archived executed notebook/dashboard\par
\textbf{Used on:} 34. 5. Operator: the record requires review\par
\textbf{Asset type:} Actual local GridGuard screenshot\par
\textbf{Transformation:} Existing HTML rendered in a browser; selected element crop; text enlarged without changing recorded values. Full provenance: research/screenshots/capture\_manifest.json.\par
\textbf{Evidence limit:} Original Explanation output. The separate acknowledgment record remains PENDING, with no dispatch or equipment operation.\par
\textbf{Original source / page:} \href{../research/source_TRACE_REST.html}{\small\ttfamily research/\allowbreak{}source\_\allowbreak{}TRACE\_\allowbreak{}REST.\allowbreak{}html}\par
\textbf{Original figure / image URL:} \href{../research/evidence/additional/rest_review.png}{\small\ttfamily research/\allowbreak{}evidence/\allowbreak{}additional/\allowbreak{}rest\_\allowbreak{}review.\allowbreak{}png}\par
\par\medskip
\subsection*{Visual credit: classroom\_cyber.png}
\textbf{Caption:} Actual L8 sweep on the four-bus fixture. The left axis uses a symmetric-log scale; the right panel is an offline audit.\par
\textbf{Credit:} Existing GridGuard workshop numerical figure\par
\textbf{Used on:} 45. A quiet residual can hide a wrong inferred state\par
\textbf{Asset type:} Recorded numerical teaching figure\par
\textbf{Transformation:} Unchanged original plot.\par
\textbf{Evidence limit:} Actual L8 sweep on the four-bus fixture. The left axis uses a symmetric-log scale; the right panel is an offline audit.\par
\textbf{Original source / page:} \href{../research/source_CY05.html}{\small\ttfamily research/\allowbreak{}source\_\allowbreak{}CY05.\allowbreak{}html}\par
\textbf{Original figure / image URL:} \href{../research/evidence/figures/classroom_cyber.png}{\small\ttfamily research/\allowbreak{}evidence/\allowbreak{}figures/\allowbreak{}classroom\_\allowbreak{}cyber.\allowbreak{}png}\par
\par\medskip
\subsection*{Visual credit: workshop\_trusted\_control.png}
\textbf{Caption:} Offline physical audit of the two decisions. Legend labels identify which estimate selected each action; bars show actual toy-model flows.\par
\textbf{Credit:} Existing GridGuard workshop numerical figure\par
\textbf{Used on:} 48. Check the consequences on every branch\par
\textbf{Asset type:} Recorded numerical teaching figure\par
\textbf{Transformation:} Unchanged original plot.\par
\textbf{Evidence limit:} Offline physical audit of the two decisions. Legend labels identify which estimate selected each action; bars show actual toy-model flows.\par
\textbf{Original source / page:} \href{../research/source_CY06.html}{\small\ttfamily research/\allowbreak{}source\_\allowbreak{}CY06.\allowbreak{}html}\par
\textbf{Original figure / image URL:} \href{../research/evidence/figures/workshop_trusted_control.png}{\small\ttfamily research/\allowbreak{}evidence/\allowbreak{}figures/\allowbreak{}workshop\_\allowbreak{}trusted\_\allowbreak{}control.\allowbreak{}png}\par
\par\medskip
\subsection*{Visual credit: cyber\_input.png}
\textbf{Caption:} Actual Awareness response containing the structured Planning input, 21 September 2026.\par
\textbf{Credit:} GridGuard project: browser capture of archived executed notebook/dashboard\par
\textbf{Used on:} 49. 1. Inputs: a declared four-bus DC task\par
\textbf{Asset type:} Actual local GridGuard screenshot\par
\textbf{Transformation:} Existing HTML rendered in a browser; selected element crop; text enlarged without changing recorded values. Full provenance: research/screenshots/capture\_manifest.json.\par
\textbf{Evidence limit:} Actual Awareness response containing the structured Planning input, 21 September 2026.\par
\textbf{Original source / page:} \href{../research/source_TRACE_CYBER.html}{\small\ttfamily research/\allowbreak{}source\_\allowbreak{}TRACE\_\allowbreak{}CYBER.\allowbreak{}html}\par
\textbf{Original figure / image URL:} \href{../research/evidence/additional/cyber_input.png}{\small\ttfamily research/\allowbreak{}evidence/\allowbreak{}additional/\allowbreak{}cyber\_\allowbreak{}input.\allowbreak{}png}\par
\par\medskip
\subsection*{Visual credit: cyber\_route.png}
\textbf{Caption:} Actual archived llama3.1:8b reply. Python executes the tool after checking the reply.\par
\textbf{Credit:} GridGuard project: browser capture of archived executed notebook/dashboard\par
\textbf{Used on:} 50. 2. Agent: choose the authorized estimator\par
\textbf{Asset type:} Actual local GridGuard screenshot\par
\textbf{Transformation:} Existing HTML rendered in a browser; selected element crop; text enlarged without changing recorded values. Full provenance: research/screenshots/capture\_manifest.json.\par
\textbf{Evidence limit:} Actual archived llama3.1:8b reply. Python executes the tool after checking the reply.\par
\textbf{Original source / page:} \href{../research/source_TRACE_CYBER.html}{\small\ttfamily research/\allowbreak{}source\_\allowbreak{}TRACE\_\allowbreak{}CYBER.\allowbreak{}html}\par
\textbf{Original figure / image URL:} \href{../research/evidence/additional/cyber_route.png}{\small\ttfamily research/\allowbreak{}evidence/\allowbreak{}additional/\allowbreak{}cyber\_\allowbreak{}route.\allowbreak{}png}\par
\par\medskip
\subsection*{Visual credit: cyber\_result.png}
\textbf{Caption:} Original injection and predicted-flow excerpt; branch order 01,12,23,03. Flow12 is the second array entry.\par
\textbf{Credit:} GridGuard project: browser capture of archived executed notebook/dashboard\par
\textbf{Used on:} 52. 4. Result: a bounded correction with evidence\par
\textbf{Asset type:} Actual local GridGuard screenshot\par
\textbf{Transformation:} Existing HTML rendered in a browser; selected element crop; text enlarged without changing recorded values. Full provenance: research/screenshots/capture\_manifest.json.\par
\textbf{Evidence limit:} Original injection and predicted-flow excerpt; branch order 01,12,23,03. Flow12 is the second array entry.\par
\textbf{Original source / page:} \href{../research/source_TRACE_CYBER.html}{\small\ttfamily research/\allowbreak{}source\_\allowbreak{}TRACE\_\allowbreak{}CYBER.\allowbreak{}html}\par
\textbf{Original figure / image URL:} \href{../research/evidence/additional/cyber_result.png}{\small\ttfamily research/\allowbreak{}evidence/\allowbreak{}additional/\allowbreak{}cyber\_\allowbreak{}result.\allowbreak{}png}\par
\par\medskip
\subsection*{Visual credit: cyber\_review.png}
\textbf{Caption:} Actual archived explanation. No human approval, dispatch authorization or equipment operation is recorded.\par
\textbf{Credit:} GridGuard project: browser capture of archived executed notebook/dashboard\par
\textbf{Used on:} 53. 5. Operator: resolve trust before proceeding\par
\textbf{Asset type:} Actual local GridGuard screenshot\par
\textbf{Transformation:} Existing HTML rendered in a browser; selected element crop; text enlarged without changing recorded values. Full provenance: research/screenshots/capture\_manifest.json.\par
\textbf{Evidence limit:} Actual archived explanation. No human approval, dispatch authorization or equipment operation is recorded.\par
\textbf{Original source / page:} \href{../research/source_TRACE_CYBER.html}{\small\ttfamily research/\allowbreak{}source\_\allowbreak{}TRACE\_\allowbreak{}CYBER.\allowbreak{}html}\par
\textbf{Original figure / image URL:} \href{../research/evidence/additional/cyber_review.png}{\small\ttfamily research/\allowbreak{}evidence/\allowbreak{}additional/\allowbreak{}cyber\_\allowbreak{}review.\allowbreak{}png}\par
\par\medskip
\end{document}
